% !TEX program = xelatex
\documentclass[UTF8,fontset=fandol,a4paper,11pt]{ctexart}
\usepackage[margin=21mm,top=23mm,bottom=22mm,headheight=16pt]{geometry}
\usepackage{amsmath,amssymb,mathtools,bm}
\usepackage{tikz}
\usetikzlibrary{arrows.meta,calc,angles,quotes,patterns}
\usepackage{xcolor,booktabs,array,longtable,enumitem,fancyhdr,needspace}
\usepackage[unicode,hidelinks]{hyperref}
\definecolor{ink}{HTML}{24343C}
\definecolor{accent}{HTML}{166A70}
\definecolor{muted}{HTML}{65747A}
\definecolor{rulegray}{HTML}{D4DDE0}
\definecolor{warm}{HTML}{966037}
\color{ink}
\setlength{\parindent}{2em}
\setlength{\parskip}{4pt}
\linespread{1.12}
\AtBeginDocument{\setlength{\abovedisplayskip}{6pt plus 2pt minus 1pt}
\setlength{\belowdisplayskip}{6pt plus 2pt minus 1pt}
\setlength{\abovedisplayshortskip}{0pt plus 1pt}
\setlength{\belowdisplayshortskip}{4pt plus 1pt}}
\setlist[enumerate]{label=(\arabic*),leftmargin=2.3em,itemsep=3pt,topsep=4pt}
\setlist[itemize]{leftmargin=2em,itemsep=2pt,topsep=4pt}
\allowdisplaybreaks[1]
\setlength{\emergencystretch}{2em}
\setcounter{secnumdepth}{0}
\providecommand{\tightlist}{\setlength{\itemsep}{0pt}\setlength{\parskip}{0pt}}
\ctexset{section={format=\Large\bfseries\color{accent},beforeskip=10pt,afterskip=9pt},
subsection={format=\normalsize\bfseries\color{accent},beforeskip=10pt,afterskip=4pt},
subsubsection={format=\normalsize\bfseries,beforeskip=8pt,afterskip=3pt}}
\pagestyle{fancy}
\fancyhf{}
\fancyhead[L]{\small\color{muted}物理学难题集萃}
% 页眉仅保留书名；版本名称只在封面显示。
\fancyfoot[L]{\small\color{muted}\chapterlabel}
\fancyfoot[R]{\small\thepage}
\renewcommand{\headrulewidth}{0.4pt}
\renewcommand{\footrulewidth}{0pt}
\newcommand{\dd}{\mathop{}\!\mathrm{d}}
\newcommand{\vect}[1]{\bm{#1}}
\newcommand{\tagline}[1]{\noindent{\small\color{muted}#1}\par}
\newcommand{\result}[1]{\begin{equation*}\boxed{#1}\end{equation*}}
\newcommand{\figcaption}[1]{\par\vspace{2pt}{\small\color{muted}#1}\par}
\newcommand{\teach}[2]{\Needspace{4\baselineskip}\subsection*{讲解要点与追问}#1\par\smallskip\noindent\textbf{追问：}#2\par}
\newcommand{\workarea}{\par\vspace{7mm}\noindent{\small\color{muted}作答区}\par\vspace{4mm}
\noindent\textcolor{rulegray}{\rule{\linewidth}{0.3pt}}\par\vspace{16mm}
\noindent\textcolor{rulegray}{\rule{\linewidth}{0.3pt}}\par\vspace{16mm}
\noindent\textcolor{rulegray}{\rule{\linewidth}{0.3pt}}\par\vspace{16mm}
\noindent\textcolor{rulegray}{\rule{\linewidth}{0.3pt}}}
\tikzset{axis/.style={-{Stealth[length=2mm]},thin,color=muted},
motion/.style={-{Stealth[length=2.3mm]},thick,color=accent},
force/.style={-{Stealth[length=2.3mm]},thick,color=warm},
traj/.style={thick,color=accent},guide/.style={dashed,thin,color=muted},
object/.style={thick,color=ink},every node/.style={font=\small}}\widowpenalty=10000\clubpenalty=10000\displaywidowpenalty=10000\newcommand{\chapterlabel}{光学第三章\quad 光的衍射}







\hypersetup{pdftitle={光学第三章 光的衍射·学生版},pdfauthor={}}

\begin{document}\pdfbookmark[0]{光学第三章 光的衍射}{cover}

\thispagestyle{empty}\vspace*{22mm}

\noindent{\Huge\bfseries 物理学难题集萃}\par\vspace{10mm}

\noindent{\LARGE 光学第三章}\par\vspace{6mm}

\noindent{\LARGE 光的衍射}\par\vspace{12mm}

\noindent{\Large\color{accent}学生版}\par\vspace{9mm}

\noindent\textcolor{accent}{\rule{\linewidth}{1pt}}\par\vspace{9mm}

\noindent 第1--22题\par\vfill\clearpage

\pdfbookmark[1]{题目目录}{contents}\section*{题目目录}

\begin{longtable}{@{}r p{113mm} r@{}}\toprule 题号 & 题目 & 页码\\\midrule

\hyperref[q:1]{1} & \hyperref[q:1]{环形孔的轴上菲涅耳衍射} & \pageref{q:1} \\

\hyperref[q:2]{2} & \hyperref[q:2]{移动屏与扩大圆孔的明暗序列} & \pageref{q:2} \\

\hyperref[q:3]{3} & \hyperref[q:3]{圆孔前后的相位与焦点增强} & \pageref{q:3} \\

\hyperref[q:4]{4} & \hyperref[q:4]{波带片的多级成像与效率} & \pageref{q:4} \\

\hyperref[q:5]{5} & \hyperref[q:5]{半缝加相位片后的衍射分布} & \pageref{q:5} \\

\hyperref[q:6]{6} & \hyperref[q:6]{余弦透光狭缝的主瓣与旁瓣} & \pageref{q:6} \\

\hyperref[q:7]{7} & \hyperref[q:7]{宽窄两缝的衍射与等效三缝} & \pageref{q:7} \\

\hyperref[q:8]{8} & \hyperref[q:8]{中缝反相的三缝衍射} & \pageref{q:8} \\

\hyperref[q:9]{9} & \hyperref[q:9]{艾里斑内的能量比例} & \pageref{q:9} \\

\hyperref[q:10]{10} & \hyperref[q:10]{环形孔的远场与中央遮挡} & \pageref{q:10} \\

\hyperref[q:11]{11} & \hyperref[q:11]{三种接收孔的孔径标度} & \pageref{q:11} \\

\hyperref[q:12]{12} & \hyperref[q:12]{有缺级要求的光栅设计} & \pageref{q:12} \\

\hyperref[q:13]{13} & \hyperref[q:13]{有限光栅的次极大位置与强度} & \pageref{q:13} \\

\hyperref[q:14]{14} & \hyperref[q:14]{显微镜的分辨距离与有效放大} & \pageref{q:14} \\

\hyperref[q:15]{15} & \hyperref[q:15]{交替间隔的成对光栅} & \pageref{q:15} \\

\hyperref[q:16]{16} & \hyperref[q:16]{中部间距误差对光栅谱峰的影响} & \pageref{q:16} \\

\hyperref[q:17]{17} & \hyperref[q:17]{两种波长的谱线重叠} & \pageref{q:17} \\

\hyperref[q:18]{18} & \hyperref[q:18]{闪耀光栅的几何与角色散} & \pageref{q:18} \\

\hyperref[q:19]{19} & \hyperref[q:19]{透射阶梯光栅的高阶光谱} & \pageref{q:19} \\

\hyperref[q:20]{20} & \hyperref[q:20]{双平面波记录的全息光栅} & \pageref{q:20} \\

\hyperref[q:21]{21} & \hyperref[q:21]{全息光栅的反衬度与衍射效率} & \pageref{q:21} \\

\hyperref[q:22]{22} & \hyperref[q:22]{点物全息图的虚像与实像} & \pageref{q:22} \\

\bottomrule\end{longtable}

\clearpage

\pdfbookmark[2]{第 1 题 环形孔的轴上菲涅耳衍射}{q1}

\section*{第 1 题\quad 环形孔的轴上菲涅耳衍射}\label{q:1}

点光源$S$发出波长$\lambda=500\,\mathrm{nm}$的单色光，距一环形孔所在平面$R=1.00\,\mathrm m$。孔内、外半径分别为$a_1=0.500\,\mathrm{mm}$和$a_2=1.00\,\mathrm{mm}$，轴上观察点$P$距孔平面$b=1.00\,\mathrm m$。求$P$点强度与不放遮板时该点强度之比。分别用半波带和衍射积分说明结果。

\begin{center}\begin{minipage}{.98\linewidth}\centering
\begin{tikzpicture}[scale=.88,transform shape]

\draw[guide] (-3.5,0)--(3.5,0);
\draw[object] (0,1.5)--(0,2);\draw[object] (0,-1.5)--(0,-2);
\draw[very thick,color=muted] (0,-.75)--(0,.75);
\fill (-3,0) circle(.035);\fill (3,0) circle(.035);\fill (0,1.125) circle(.03);
% paths:Q-upper and Q-lower; transverse 1.5 cm/mm, axial 3 cm/m
\draw[motion] (-3,0)--(0,1.125);\draw[motion] (0,1.125)--(3,0);
\draw[motion] (-3,0)--(0,-1.125);\draw[motion] (0,-1.125)--(3,0);
\node[below]at(-3,0){$S$};\node[right]at(3,0){$P$};
\node[below right]at(0,0){$O$};\node[above left,xshift=-2pt]at(0,1.125){$Q$};
\draw[<->,thin,color=muted] (-3,-2.35)--(0,-2.35);\node[below]at(-1.5,-2.35){$R=1.00\,\mathrm m$};
\draw[<->,thin,color=muted] (0,-2.35)--(3,-2.35);\node[below]at(1.5,-2.35){$b=1.00\,\mathrm m$};
\draw[<->,thin,color=muted] (.6,0)--(.6,1.5);\node[right,fill=white,inner sep=1pt]at(.6,.68){$a_2$};
\draw[<->,thin,color=muted] (-.65,0)--(-.65,.75);\node[left,fill=white,inner sep=1pt]at(-.65,.35){$a_1$};
\node at(0,2.35){纵向毫米、轴向米，异尺度；两段为衍射代表路径};

\end{tikzpicture}
\figcaption{图光3.1\quad 等距点源与观察点，中心及外侧不透光}
\end{minipage}\end{center}

\workarea

\clearpage

\pdfbookmark[2]{第 2 题 移动屏与扩大圆孔的明暗序列}{q2}

\section*{第 2 题\quad 移动屏与扩大圆孔的明暗序列}\label{q:2}

波长$\lambda=500\,\mathrm{nm}$的点光源距半径$a=2.00\,\mathrm{mm}$的圆孔$R=2.00\,\mathrm m$。（1）观察屏由很远处向圆孔移动，求中心依次出现的前三个亮位置。（2）求前三个有限距离的暗位置。（3）把屏固定在孔后$b=5.00\,\mathrm m$，从$a=2.00\,\mathrm{mm}$起逐渐增大孔半径，求首次变暗和首次变亮时的半径。采用菲涅耳近轴近似。
亮、暗按相对于同一位置无遮挡强度的衍射因子判断。
\begin{center}\begin{minipage}{.98\linewidth}\centering
\begin{tikzpicture}[scale=.85,transform shape]
\draw[axis](-.4,0)--(7.5,0);\fill(0,0)circle(.035);\node[below]at(0,0){$S$};
\fill(2,0)circle(.025);\node[below left]at(2,0){$O$};\fill(7,0)circle(.035);\node[above right,xshift=2pt,yshift=3pt]at(7,0){$P$};
\draw[object](2,1.1)--(2,1.65)(2,-1.1)--(2,-1.65);
\draw[object](7,-1.6)--(7,1.6);\node[above]at(7,1.65){屏};
\draw[motion](0,0)--(2,1.1);\draw[motion](2,1.1)--(7,0);
\draw[guide](0,0)--(2,-1.1)--(7,0);\fill(2,1.1)circle(.03);\node[above right]at(2,1.1){$Q$};
\draw[<->](0,-1.85)--(2,-1.85);\node[below]at(1,-1.85){$R=2\,\mathrm m$};
\draw[<->](2,-1.85)--(7,-1.85);\node[below]at(4.5,-1.85){$b$};
\draw[<->](2.55,0)--(2.55,1.1);\node[right]at(2.55,.55){$a=2\,\mathrm{mm}$};
\node[above]at(3.7,2){横向距离按米；纵向尺寸放大};
\end{tikzpicture}
\figcaption{图光3.2\quad 点光源、圆孔与移动屏的轴截面}
\end{minipage}\end{center}

\workarea

\clearpage

\pdfbookmark[2]{第 3 题 圆孔前后的相位与焦点增强}{q3}

\section*{第 3 题\quad 圆孔前后的相位与焦点增强}\label{q:3}

强度为$I_0$、波长$\lambda=550\,\mathrm{nm}$的单色平行光垂直入射直径$d=1.10\,\mathrm{mm}$的圆孔。轴上点$P$在孔后$b=33.0\,\mathrm{cm}$。（1）求$P$点强度。（2）在孔内放一片无吸收的理想薄会聚透镜，焦距$f=33.0\,\mathrm{cm}$，再求$P$点强度。用相位矢量说明两种结果的差别。

\begin{center}\begin{minipage}{.98\linewidth}\centering
\begin{tikzpicture}[x=.022cm,y=1.7cm,every node/.style={font=\small}]
\draw[axis](-70,0)--(365,0) node[right]{$z$};\draw[object](0,-.55)--(0,-.85);\draw[object](0,.55)--(0,.85);\draw[object,<->](0,-.55)--(0,.55);
\draw[motion](-60.000000000000,-0.400000000000)--(0.000000000000,-0.400000000000);\draw[guide](0.000000000000,-0.400000000000)--(330,0);
\draw[motion](-60.000000000000,0.000000000000)--(0.000000000000,0.000000000000);\draw[guide](0.000000000000,0.000000000000)--(330,0);
\draw[motion](-60.000000000000,0.400000000000)--(0.000000000000,0.400000000000);\draw[guide](0.000000000000,0.400000000000)--(330,0);
\fill(330,0)circle(1.2pt);\node[above right]at(330.000000000000,0.000000000000){$P$};
\draw[<->,muted](-80,-.55)--(-80,.55) node[midway,left]{$d=1.10\,\mathrm{mm}$};\draw[<->,muted](0,-1)--(330,-1) node[midway,below]{$b=330\,\mathrm{mm}$};\node[above]at(0.000000000000,0.900000000000){孔面，可置透镜};
\node[above]at(190.000000000000,0.650000000000){平行入射；纵向放大};
\end{tikzpicture}
\figcaption{图光3.3\quad 圆孔、轴上点与可置于孔面的透镜；虚线为积分传播路径}
\end{minipage}\end{center}

\workarea

\clearpage

\pdfbookmark[2]{第 4 题 波带片的多级成像与效率}{q4}

\section*{第 4 题\quad 波带片的多级成像与效率}\label{q:4}

菲涅耳波带片共有$2N$个半波带，奇数带透明、偶数带不透明，中心透明带直径$d=1.00\,\mathrm{mm}$。以波长$\lambda=500\,\mathrm{nm}$的光照明，轴上点光源位于片前$s=1.00\,\mathrm m$。（1）求波带片的实像位置。（2）改用相同外孔径、焦距等于波带片主焦距的无吸收薄会聚透镜，比较主像点的强度。

\begin{center}\begin{minipage}{.98\linewidth}\centering
\begin{tikzpicture}[scale=0.88,transform shape]
\fill[fill=ink!25](-4.000000000000,0.000000000000)circle(1.697056274848);\draw[thin](-4.000000000000,0.000000000000)circle(1.697056274848);\fill[fill=white](-4.000000000000,0.000000000000)circle(1.587450786639);\draw[thin](-4.000000000000,0.000000000000)circle(1.587450786639);\fill[fill=ink!25](-4.000000000000,0.000000000000)circle(1.469693845670);\draw[thin](-4.000000000000,0.000000000000)circle(1.469693845670);\fill[fill=white](-4.000000000000,0.000000000000)circle(1.341640786500);\draw[thin](-4.000000000000,0.000000000000)circle(1.341640786500);\fill[fill=ink!25](-4.000000000000,0.000000000000)circle(1.200000000000);\draw[thin](-4.000000000000,0.000000000000)circle(1.200000000000);\fill[fill=white](-4.000000000000,0.000000000000)circle(1.039230484541);\draw[thin](-4.000000000000,0.000000000000)circle(1.039230484541);\fill[fill=ink!25](-4.000000000000,0.000000000000)circle(0.848528137424);\draw[thin](-4.000000000000,0.000000000000)circle(0.848528137424);\fill[fill=white](-4.000000000000,0.000000000000)circle(0.600000000000);\draw[thin](-4.000000000000,0.000000000000)circle(0.600000000000);\node[above]at(-4.000000000000,2.050000000000){$N=4,\quad r_m=\sqrt m\,a_1$};\draw[<->,thin](-4.600000000000,-2.000000000000)--(-3.400000000000,-2.000000000000);\node[below]at(-4.000000000000,-2.000000000000){$d=2a_1=1\,\mathrm{mm}$};\draw[axis](-0.200000000000,0.000000000000)--(6.400000000000,0.000000000000);\draw[guide](3.000000000000,-1.900000000000)--(3.000000000000,1.900000000000);\draw[object](3.000000000000,-0.600000000000)--(3.000000000000,-0.848528137424);\draw[object](3.000000000000,0.600000000000)--(3.000000000000,0.848528137424);\draw[object](3.000000000000,-1.039230484541)--(3.000000000000,-1.200000000000);\draw[object](3.000000000000,1.039230484541)--(3.000000000000,1.200000000000);\draw[object](3.000000000000,-1.341640786500)--(3.000000000000,-1.469693845670);\draw[object](3.000000000000,1.341640786500)--(3.000000000000,1.469693845670);\draw[object](3.000000000000,-1.587450786639)--(3.000000000000,-1.697056274848);\draw[object](3.000000000000,1.587450786639)--(3.000000000000,1.697056274848);\fill[color=ink](0.000000000000,0.000000000000)circle(1.15pt);\fill[color=ink](6.000000000000,0.000000000000)circle(1.15pt);\node[below]at(0.000000000000,0.000000000000){$S$};\node[below left]at(3.000000000000,-0.020000000000){$O$};\node[above right,yshift=2pt]at(6.000000000000,0.000000000000){$P_1$};\draw[motion](0.000000000000,0.000000000000)--(3.000000000000,1.550000000000);\draw[motion](3.000000000000,1.550000000000)--(6.000000000000,0.000000000000);\node[above left]at(3.000000000000,1.550000000000){$Q$};\draw[<->,thin](0.000000000000,-2.050000000000)--(3.000000000000,-2.050000000000);\draw[<->,thin](3.000000000000,-2.050000000000)--(6.000000000000,-2.050000000000);\node[below]at(1.500000000000,-2.050000000000){$s=1\,\mathrm m$};\node[below]at(4.500000000000,-2.050000000000){$b_1=1\,\mathrm m$};\node[above]at(3.000000000000,2.200000000000){纵向孔径用mm、轴向用m；异尺度};
\end{tikzpicture}
\figcaption{图光3.4\quad 八带波带片的真实平方根边界与主像几何}\end{minipage}\end{center}

\workarea

\clearpage

\pdfbookmark[2]{第 5 题 半缝加相位片后的衍射分布}{q5}

\section*{第 5 题\quad 半缝加相位片后的衍射分布}\label{q:5}

宽度为$a$的狭缝，一半缝前覆盖理想透明相位片，使该半缝的透射光多出$\pi$相位，振幅不变。波长为$\lambda$的平行光正入射。（1）求夫琅禾费衍射强度分布，以未加片时同一狭缝的中央强度$I_c$归一化。（2）取$a=10\lambda$，画出分布，与普通单缝比较，并说明极小及最近极大的位置。

\begin{center}\begin{minipage}{.98\linewidth}\centering
\begin{tikzpicture}[scale=1.1]

\draw[axis](-2,0)--(4.1,0)node[right]{$z$};
\draw[axis](0,-1.6)--(0,1.7)node[above]{$x$};
\draw[object](0,-1.5)--(0,-1);\draw[object](0,1)--(0,1.5);
\draw[very thick,color=warm](-.18,0)--(-.18,1);
\node[above left]at(-.18,1.02){$t=-1$};
\node[below left]at(-1.3,-1.2){$t=+1$};
\draw[motion](-1.8,.5)--(-.22,.5);\draw[motion](-1.8,-.5)--(0,-.5);
\fill(0,.5)circle(.025);\fill(0,-.5)circle(.025);
\node[above right]at(.05,.5){$x=a/4$};
\node[below left]at(-.05,-.5){$x=-a/4$};
\draw[motion](0,.5)--(2.814582562,-1.125);
\draw[motion](0,-.5)--(2.381569860,-1.875);
\draw[guide](2.208364780,-2.175)--(2.987787642,-.825);
\node[right]at(3.02,-.75){$\Sigma$};
\draw[thin](1,0)arc[start angle=0,end angle=-30,radius=1];
\node[above]at(1.3,.12){$\theta>0$};
\draw[<->,thin](3.55,-1)--(3.55,1);
\node[right]at(3.55,.4){$a$};
\node[below]at(1,-2.35){指定远场方向；$\Sigma$与出射方向垂直};
\end{tikzpicture}
\figcaption{图光3.5}
\end{minipage}\end{center}

\workarea

\clearpage

\pdfbookmark[2]{第 6 题 余弦透光狭缝的主瓣与旁瓣}{q6}

\section*{第 6 题\quad 余弦透光狭缝的主瓣与旁瓣}\label{q:6}

宽度为$a$的狭缝具有振幅透过率$t(x)=\cos(\pi x/a)$，$-a/2\le x\le a/2$，缝外为零。波长$\lambda$的单色平行光正入射。（1）求夫琅禾费强度分布及极小位置。（2）计算相邻正侧极小之间，在$\sin\theta$坐标的中点处强度与中央强度之比，比较普通单缝，说明这种渐变透光的作用。

\begin{center}\begin{minipage}{.98\linewidth}\centering
\begin{tikzpicture}[scale=.95,transform shape]
\draw[axis] (0.00000,0.00000)--(5.25000,0.00000);\draw[axis] (0.00000,0.00000)--(0.00000,2.05000);\node[right] at (5.25000,0.00000) {$x/a$};\node[above right] at (0.00000,2.05000) {$t$};\draw[traj] plot coordinates {(0.00000,0.00000) (0.01000,0.01131) (0.02000,0.02262) (0.03000,0.03393) (0.04000,0.04523) (0.05000,0.05654) (0.06000,0.06784) (0.07000,0.07914) (0.08000,0.09044) (0.09000,0.10173) (0.10000,0.11302) (0.11000,0.12431) (0.12000,0.13559) (0.13000,0.14686) (0.14000,0.15813) (0.15000,0.16939) (0.16000,0.18065) (0.17000,0.19190) (0.18000,0.20314) (0.19000,0.21437) (0.20000,0.22560) (0.21000,0.23682) (0.22000,0.24802) (0.23000,0.25922) (0.24000,0.27041) (0.25000,0.28158) (0.26000,0.29275) (0.27000,0.30390) (0.28000,0.31504) (0.29000,0.32617) (0.30000,0.33729) (0.31000,0.34839) (0.32000,0.35948) (0.33000,0.37055) (0.34000,0.38161) (0.35000,0.39266) (0.36000,0.40369) (0.37000,0.41470) (0.38000,0.42570) (0.39000,0.43668) (0.40000,0.44764) (0.41000,0.45859) (0.42000,0.46951) (0.43000,0.48042) (0.44000,0.49131) (0.45000,0.50218) (0.46000,0.51303) (0.47000,0.52387) (0.48000,0.53467) (0.49000,0.54546) (0.50000,0.55623) (0.51000,0.56698) (0.52000,0.57770) (0.53000,0.58840) (0.54000,0.59908) (0.55000,0.60973) (0.56000,0.62036) (0.57000,0.63096) (0.58000,0.64154) (0.59000,0.65210) (0.60000,0.66262) (0.61000,0.67313) (0.62000,0.68360) (0.63000,0.69405) (0.64000,0.70447) (0.65000,0.71487) (0.66000,0.72523) (0.67000,0.73557) (0.68000,0.74588) (0.69000,0.75615) (0.70000,0.76640) (0.71000,0.77662) (0.72000,0.78681) (0.73000,0.79696) (0.74000,0.80709) (0.75000,0.81718) (0.76000,0.82724) (0.77000,0.83727) (0.78000,0.84727) (0.79000,0.85723) (0.80000,0.86716) (0.81000,0.87705) (0.82000,0.88691) (0.83000,0.89673) (0.84000,0.90652) (0.85000,0.91627) (0.86000,0.92599) (0.87000,0.93567) (0.88000,0.94531) (0.89000,0.95492) (0.90000,0.96449) (0.91000,0.97402) (0.92000,0.98351) (0.93000,0.99296) (0.94000,1.00238) (0.95000,1.01175) (0.96000,1.02108) (0.97000,1.03038) (0.98000,1.03963) (0.99000,1.04884) (1.00000,1.05801) (1.01000,1.06714) (1.02000,1.07623) (1.03000,1.08527) (1.04000,1.09427) (1.05000,1.10323) (1.06000,1.11215) (1.07000,1.12102) (1.08000,1.12984) (1.09000,1.13863) (1.10000,1.14736) (1.11000,1.15605) (1.12000,1.16470) (1.13000,1.17330) (1.14000,1.18185) (1.15000,1.19036) (1.16000,1.19882) (1.17000,1.20723) (1.18000,1.21560) (1.19000,1.22392) (1.20000,1.23218) (1.21000,1.24040) (1.22000,1.24858) (1.23000,1.25670) (1.24000,1.26477) (1.25000,1.27279) (1.26000,1.28076) (1.27000,1.28869) (1.28000,1.29656) (1.29000,1.30438) (1.30000,1.31214) (1.31000,1.31986) (1.32000,1.32752) (1.33000,1.33514) (1.34000,1.34269) (1.35000,1.35020) (1.36000,1.35765) (1.37000,1.36505) (1.38000,1.37240) (1.39000,1.37969) (1.40000,1.38692) (1.41000,1.39411) (1.42000,1.40123) (1.43000,1.40830) (1.44000,1.41532) (1.45000,1.42228) (1.46000,1.42918) (1.47000,1.43603) (1.48000,1.44282) (1.49000,1.44955) (1.50000,1.45623) (1.51000,1.46285) (1.52000,1.46941) (1.53000,1.47591) (1.54000,1.48236) (1.55000,1.48875) (1.56000,1.49507) (1.57000,1.50134) (1.58000,1.50755) (1.59000,1.51370) (1.60000,1.51979) (1.61000,1.52582) (1.62000,1.53179) (1.63000,1.53770) (1.64000,1.54355) (1.65000,1.54934) (1.66000,1.55506) (1.67000,1.56073) (1.68000,1.56633) (1.69000,1.57187) (1.70000,1.57735) (1.71000,1.58277) (1.72000,1.58812) (1.73000,1.59342) (1.74000,1.59865) (1.75000,1.60381) (1.76000,1.60891) (1.77000,1.61395) (1.78000,1.61893) (1.79000,1.62384) (1.80000,1.62869) (1.81000,1.63347) (1.82000,1.63819) (1.83000,1.64284) (1.84000,1.64743) (1.85000,1.65196) (1.86000,1.65642) (1.87000,1.66081) (1.88000,1.66514) (1.89000,1.66940) (1.90000,1.67360) (1.91000,1.67773) (1.92000,1.68179) (1.93000,1.68579) (1.94000,1.68972) (1.95000,1.69359) (1.96000,1.69738) (1.97000,1.70111) (1.98000,1.70478) (1.99000,1.70837) (2.00000,1.71190) (2.01000,1.71536) (2.02000,1.71876) (2.03000,1.72208) (2.04000,1.72534) (2.05000,1.72853) (2.06000,1.73165) (2.07000,1.73470) (2.08000,1.73769) (2.09000,1.74060) (2.10000,1.74345) (2.11000,1.74623) (2.12000,1.74894) (2.13000,1.75158) (2.14000,1.75415) (2.15000,1.75665) (2.16000,1.75908) (2.17000,1.76145) (2.18000,1.76374) (2.19000,1.76596) (2.20000,1.76812) (2.21000,1.77020) (2.22000,1.77222) (2.23000,1.77416) (2.24000,1.77603) (2.25000,1.77784) (2.26000,1.77957) (2.27000,1.78124) (2.28000,1.78283) (2.29000,1.78435) (2.30000,1.78581) (2.31000,1.78719) (2.32000,1.78850) (2.33000,1.78974) (2.34000,1.79091) (2.35000,1.79201) (2.36000,1.79304) (2.37000,1.79400) (2.38000,1.79489) (2.39000,1.79570) (2.40000,1.79645) (2.41000,1.79712) (2.42000,1.79773) (2.43000,1.79826) (2.44000,1.79872) (2.45000,1.79911) (2.46000,1.79943) (2.47000,1.79968) (2.48000,1.79986) (2.49000,1.79996) (2.50000,1.80000) (2.51000,1.79996) (2.52000,1.79986) (2.53000,1.79968) (2.54000,1.79943) (2.55000,1.79911) (2.56000,1.79872) (2.57000,1.79826) (2.58000,1.79773) (2.59000,1.79712) (2.60000,1.79645) (2.61000,1.79570) (2.62000,1.79489) (2.63000,1.79400) (2.64000,1.79304) (2.65000,1.79201) (2.66000,1.79091) (2.67000,1.78974) (2.68000,1.78850) (2.69000,1.78719) (2.70000,1.78581) (2.71000,1.78435) (2.72000,1.78283) (2.73000,1.78124) (2.74000,1.77957) (2.75000,1.77784) (2.76000,1.77603) (2.77000,1.77416) (2.78000,1.77222) (2.79000,1.77020) (2.80000,1.76812) (2.81000,1.76596) (2.82000,1.76374) (2.83000,1.76145) (2.84000,1.75908) (2.85000,1.75665) (2.86000,1.75415) (2.87000,1.75158) (2.88000,1.74894) (2.89000,1.74623) (2.90000,1.74345) (2.91000,1.74060) (2.92000,1.73769) (2.93000,1.73470) (2.94000,1.73165) (2.95000,1.72853) (2.96000,1.72534) (2.97000,1.72208) (2.98000,1.71876) (2.99000,1.71536) (3.00000,1.71190) (3.01000,1.70837) (3.02000,1.70478) (3.03000,1.70111) (3.04000,1.69738) (3.05000,1.69359) (3.06000,1.68972) (3.07000,1.68579) (3.08000,1.68179) (3.09000,1.67773) (3.10000,1.67360) (3.11000,1.66940) (3.12000,1.66514) (3.13000,1.66081) (3.14000,1.65642) (3.15000,1.65196) (3.16000,1.64743) (3.17000,1.64284) (3.18000,1.63819) (3.19000,1.63347) (3.20000,1.62869) (3.21000,1.62384) (3.22000,1.61893) (3.23000,1.61395) (3.24000,1.60891) (3.25000,1.60381) (3.26000,1.59865) (3.27000,1.59342) (3.28000,1.58812) (3.29000,1.58277) (3.30000,1.57735) (3.31000,1.57187) (3.32000,1.56633) (3.33000,1.56073) (3.34000,1.55506) (3.35000,1.54934) (3.36000,1.54355) (3.37000,1.53770) (3.38000,1.53179) (3.39000,1.52582) (3.40000,1.51979) (3.41000,1.51370) (3.42000,1.50755) (3.43000,1.50134) (3.44000,1.49507) (3.45000,1.48875) (3.46000,1.48236) (3.47000,1.47591) (3.48000,1.46941) (3.49000,1.46285) (3.50000,1.45623) (3.51000,1.44955) (3.52000,1.44282) (3.53000,1.43603) (3.54000,1.42918) (3.55000,1.42228) (3.56000,1.41532) (3.57000,1.40830) (3.58000,1.40123) (3.59000,1.39411) (3.60000,1.38692) (3.61000,1.37969) (3.62000,1.37240) (3.63000,1.36505) (3.64000,1.35765) (3.65000,1.35020) (3.66000,1.34269) (3.67000,1.33514) (3.68000,1.32752) (3.69000,1.31986) (3.70000,1.31214) (3.71000,1.30438) (3.72000,1.29656) (3.73000,1.28869) (3.74000,1.28076) (3.75000,1.27279) (3.76000,1.26477) (3.77000,1.25670) (3.78000,1.24858) (3.79000,1.24040) (3.80000,1.23218) (3.81000,1.22392) (3.82000,1.21560) (3.83000,1.20723) (3.84000,1.19882) (3.85000,1.19036) (3.86000,1.18185) (3.87000,1.17330) (3.88000,1.16470) (3.89000,1.15605) (3.90000,1.14736) (3.91000,1.13863) (3.92000,1.12984) (3.93000,1.12102) (3.94000,1.11215) (3.95000,1.10323) (3.96000,1.09427) (3.97000,1.08527) (3.98000,1.07623) (3.99000,1.06714) (4.00000,1.05801) (4.01000,1.04884) (4.02000,1.03963) (4.03000,1.03038) (4.04000,1.02108) (4.05000,1.01175) (4.06000,1.00238) (4.07000,0.99296) (4.08000,0.98351) (4.09000,0.97402) (4.10000,0.96449) (4.11000,0.95492) (4.12000,0.94531) (4.13000,0.93567) (4.14000,0.92599) (4.15000,0.91627) (4.16000,0.90652) (4.17000,0.89673) (4.18000,0.88691) (4.19000,0.87705) (4.20000,0.86716) (4.21000,0.85723) (4.22000,0.84727) (4.23000,0.83727) (4.24000,0.82724) (4.25000,0.81718) (4.26000,0.80709) (4.27000,0.79696) (4.28000,0.78681) (4.29000,0.77662) (4.30000,0.76640) (4.31000,0.75615) (4.32000,0.74588) (4.33000,0.73557) (4.34000,0.72523) (4.35000,0.71487) (4.36000,0.70447) (4.37000,0.69405) (4.38000,0.68360) (4.39000,0.67313) (4.40000,0.66262) (4.41000,0.65210) (4.42000,0.64154) (4.43000,0.63096) (4.44000,0.62036) (4.45000,0.60973) (4.46000,0.59908) (4.47000,0.58840) (4.48000,0.57770) (4.49000,0.56698) (4.50000,0.55623) (4.51000,0.54546) (4.52000,0.53467) (4.53000,0.52387) (4.54000,0.51303) (4.55000,0.50218) (4.56000,0.49131) (4.57000,0.48042) (4.58000,0.46951) (4.59000,0.45859) (4.60000,0.44764) (4.61000,0.43668) (4.62000,0.42570) (4.63000,0.41470) (4.64000,0.40369) (4.65000,0.39266) (4.66000,0.38161) (4.67000,0.37055) (4.68000,0.35948) (4.69000,0.34839) (4.70000,0.33729) (4.71000,0.32617) (4.72000,0.31504) (4.73000,0.30390) (4.74000,0.29275) (4.75000,0.28158) (4.76000,0.27041) (4.77000,0.25922) (4.78000,0.24802) (4.79000,0.23682) (4.80000,0.22560) (4.81000,0.21437) (4.82000,0.20314) (4.83000,0.19190) (4.84000,0.18065) (4.85000,0.16939) (4.86000,0.15813) (4.87000,0.14686) (4.88000,0.13559) (4.89000,0.12431) (4.90000,0.11302) (4.91000,0.10173) (4.92000,0.09044) (4.93000,0.07914) (4.94000,0.06784) (4.95000,0.05654) (4.96000,0.04523) (4.97000,0.03393) (4.98000,0.02262) (4.99000,0.01131) (5.00000,0.00000)};\draw[thin] (0.00000,-0.06000)--(0.00000,0.06000);\node[below] at (0.00000,-0.13000) {$-1/2$};\draw[thin] (2.50000,-0.06000)--(2.50000,0.06000);\node[below] at (2.50000,-0.13000) {0};\draw[thin] (5.00000,-0.06000)--(5.00000,0.06000);\node[below] at (5.00000,-0.13000) {$1/2$};\node[above] at (2.50000,2.50000) {振幅透过率};
\end{tikzpicture}
\figcaption{图光3.6\quad 余弦加权孔径在边缘连续归零}
\end{minipage}\end{center}

\workarea

\clearpage

\pdfbookmark[2]{第 7 题 宽窄两缝的衍射与等效三缝}{q7}

\section*{第 7 题\quad 宽窄两缝的衍射与等效三缝}\label{q:7}

两条平行狭缝的宽度分别为$a_1,a_2$，中心间距为$d$，波长为$\lambda$的平行光正入射。（1）求夫琅禾费强度分布。（2）当$a_1=2a_2=2a$、$d=3a$时，求简化表达式，并用把宽缝分成两条相邻窄缝的方法再推导一次，讨论图样的特征。

\begin{center}\begin{minipage}{.98\linewidth}\centering
\begin{tikzpicture}[x=.13cm,y=.65cm,every node/.style={font=\small}]
% GEOMETRY Q: x is downward; a=1 unit, lens z=8, f=60, screen z=68, input diffraction slope=.02.
\draw[axis](-14,0)--(73,0)node[right]{$z$};
\draw[object](0,2)--(0,1);\draw[object](0,-1)--(0,-2.5);\draw[object](0,-3.5)--(0,-4.1);
\draw[very thick,color=accent](0,-1)--(0,1);\draw[very thick,color=accent](0,-3.5)--(0,-2.5);
\node[above left]at(-.5,1){$a_1=2a$};\node[above left]at(-8,-2.5){$a_2=a$};
\draw[<->,thin,color=muted](-5,0)--(-5,-3);\node[left]at(-5,-1.6){$d=3a$};
\draw[motion](-14,0)--(-1,0);\draw[motion](-14,-3)--(-1,-3);
\draw[motion](0,0)--(8,.16);\draw[motion](0,-3)--(8,-2.84);
\draw[motion](8,.16)--(68,1.2);\draw[motion](8,-2.84)--(68,1.2);
\draw[object](8,-3.5)--(8,2);\draw[object](7,-3.25)--(8,-3.5)--(9,-3.25);\draw[object](7,1.75)--(8,2)--(9,1.75);\node[above]at(8,2.1){$L$};
\draw[object](68,-3.6)--(68,2.2);\node[above]at(68,2.3){后焦面};
\fill(68,1.2)circle(1pt);\node[right]at(69,1.2){$P$};\node[below right]at(68,-.05){$O$};
\node[above]at(30,.65){$\theta>0$向上};
\draw[axis](1,-1.4)--(1,-2.25)node[right]{$x$};
\node[below]at(35,-4.2){近轴；纵向放大；正$x$指向窄缝};
\end{tikzpicture}
\figcaption{图光3.7\quad 特殊尺寸与平行方向衍射光的后焦面会聚}
\end{minipage}\end{center}

\workarea

\clearpage

\pdfbookmark[2]{第 8 题 中缝反相的三缝衍射}{q8}

\section*{第 8 题\quad 中缝反相的三缝衍射}\label{q:8}

三条等宽$a$的平行狭缝，相邻中心距均为$d$。中间缝前放一透明相位片，使其中透射光增加$\pi$相位而振幅不变。波长为$\lambda$的平行光正入射，在透镜后焦面观察。（1）推导强度分布。（2）求第一级单缝衍射极小、最近的强干涉极大及最近的干涉极小的角位置，并区分纯干涉峰与完整强度峰。

\begin{center}\begin{minipage}{.98\linewidth}\centering
\begin{tikzpicture}[scale=.8,transform shape]
\draw[axis] (-2.50000,0.00000)--(6.80000,0.00000);\draw[object] (0.00000,-2.40000)--(0.00000,3.20000);\draw[object] (6.00000,-2.40000)--(6.00000,3.20000);\node[above] at (6.00000,3.20000) {焦平面};\draw[very thick,color=white] (0.00000,1.20000)--(0.00000,1.60000);\draw[thin] (-0.16000,1.20000)--(0.16000,1.20000);\draw[thin] (-0.16000,1.60000)--(0.16000,1.60000);\draw[motion] (-2.20000,1.40000)--(0.00000,1.40000);
\draw[motion] (0.00000,1.40000)--(3.00000,2.60000);
\draw[motion] (3.00000,2.60000)--(6.00000,1.20000);\node[above left] at (-0.12000,1.50000) {$1$};\draw[very thick,color=white] (0.00000,-0.20000)--(0.00000,0.20000);\draw[thin] (-0.16000,-0.20000)--(0.16000,-0.20000);\draw[thin] (-0.16000,0.20000)--(0.16000,0.20000);\draw[motion] (-2.20000,0.00000)--(0.00000,0.00000);
\draw[motion] (0.00000,0.00000)--(3.00000,1.20000);
\draw[motion] (3.00000,1.20000)--(6.00000,1.20000);\node[above left] at (-0.12000,0.10000) {$2$};\draw[very thick,color=accent] (-0.45000,-0.20000)--(-0.45000,0.20000);\node[below] at (-0.80000,-0.35000) {$\pi$};\draw[very thick,color=white] (0.00000,-1.60000)--(0.00000,-1.20000);\draw[thin] (-0.16000,-1.60000)--(0.16000,-1.60000);\draw[thin] (-0.16000,-1.20000)--(0.16000,-1.20000);\draw[motion] (-2.20000,-1.40000)--(0.00000,-1.40000);
\draw[motion] (0.00000,-1.40000)--(3.00000,-0.20000);
\draw[motion] (3.00000,-0.20000)--(6.00000,1.20000);\node[above left] at (-0.12000,-1.30000) {$3$};\fill (6.00000,1.20000) circle(.035);\node[right] at (6.00000,1.20000) {$P$};\fill (6.00000,0.00000) circle(.035);\node[right] at (6.00000,0.00000) {$O$};\node[above] at (2.00000,0.35000) {$\theta$};\draw[object] (3.00000,-3.00000)--(3.00000,3.00000);\draw[object] (2.87000,-2.80000)--(3.00000,-3.00000)--(3.13000,-2.80000);\draw[object] (2.87000,2.80000)--(3.00000,3.00000)--(3.13000,2.80000);\node[above] at (3.00000,3.10000) {$L$};

\draw[axis] (-3.80000,1.90000)--(-3.80000,-1.90000) node[below] {$x$};
\draw[thin] (-3.90000,1.40000)--(-3.70000,1.40000);\node[left] at (-3.95000,1.40000) {$0$};
\draw[thin] (-3.90000,0.00000)--(-3.70000,0.00000);\node[left] at (-3.95000,0.00000) {$d$};
\draw[thin] (-3.90000,-1.40000)--(-3.70000,-1.40000);\node[left] at (-3.95000,-1.40000) {$2d$};
\draw[<->,thin,color=muted] (-2.90000,1.40000)--(-2.90000,0.00000);\node[left] at (-3.00000,0.70000) {$d$};
\draw[<->,thin,color=muted] (-2.90000,0.00000)--(-2.90000,-1.40000);\node[left] at (-3.00000,-0.70000) {$d$};
\node[align=center] at (-1.60000,2.65000) {各缝宽均为$a$};
\end{tikzpicture}
\figcaption{图光3.8\quad 中缝相位片与三缝的等间距几何}
\end{minipage}\end{center}

\workarea

\clearpage

\pdfbookmark[2]{第 9 题 艾里斑内的能量比例}{q9}

\section*{第 9 题\quad 艾里斑内的能量比例}\label{q:9}

半径为$a$的圆孔以平行单色光正入射，透镜焦距为$f$，焦平面强度为$I(r)=I_c[2J_1(x)/x]^2$，其中$x=2\pi ar/(\lambda f)$，$I_c=I_{\rm in}(\pi a^2)^2/(\lambda^2f^2)$。求半径$r_0$内的能量与入射到圆孔总能量之比，并计算第一暗环以内艾里斑的能量比例。
其中$J_1$为一阶贝塞尔函数，$r=OP$，入射强度$I_{\rm in}$均匀。衍射角为$\theta$时$x=(2\pi a/\lambda)\sin\theta$；在小角条件下取$r\simeq f\theta$，并假设透镜不另行截光。

\begin{center}\begin{minipage}{.98\linewidth}\centering
\begin{tikzpicture}[x=.3cm,y=.75cm,every node/.style={font=\small}]
% GEOMETRY: a=1, lens z=4, f=20, focal plane z=24; common incoming diffraction slope .06.
\draw[axis](-7,0)--(27,0)node[right]{$z$};
\draw[object](0,1)--(0,1.8);\draw[object](0,-1)--(0,-1.8);\draw[very thick,color=accent](0,-1)--(0,1);
\draw[<->,thin,color=muted](-3,0)--(-3,1);\node[left]at(-3,.8){$a$};
\draw[motion](-7,.5)--(-.5,.5);\draw[motion](-7,-.5)--(-.5,-.5);
\draw[motion](0,.5)--(4,.74);\draw[motion](0,-.5)--(4,-.26);
\draw[motion](4,.74)--(24,1.2);\draw[motion](4,-.26)--(24,1.2);
\draw[object](4,-1.6)--(4,1.7);\draw[object](3.2,1.5)--(4,1.7)--(4.8,1.5);\draw[object](3.2,-1.4)--(4,-1.6)--(4.8,-1.4);\node[above]at(4,1.85){$L$};
\draw[object](24,-1.7)--(24,1.9);\node[above]at(24,2){后焦面};\fill(24,1.2)circle(1.2pt);\node[right]at(24,1.2){$P$};\node[below right]at(24,0){$O$};
\draw[<->,thin,color=muted](26,0)--(26,1.2);\node[right]at(26,.6){$r$};\node[above]at(13,1.4){同一衍射方向};
\draw[<->,thin,color=muted](4,-2.1)--(24,-2.1);\node[below]at(14,-2.1){$f$};\node[below]at(9,-3){近轴；纵向放大；圆孔与接收区均为圆盘};
\end{tikzpicture}
\figcaption{图光3.9\quad 圆孔截面及同方向分量的会聚}
\end{minipage}\end{center}

\workarea

\clearpage

\pdfbookmark[2]{第 10 题 环形孔的远场与中央遮挡}{q10}

\section*{第 10 题\quad 环形孔的远场与中央遮挡}\label{q:10}

外半径$a$、内半径$b$的通光圆环，受波长$\lambda$的单色平行光正入射。（1）求夫琅禾费强度分布。（2）$a=2b$时，求其中央强度与同样入射场下半径$a$的完整圆孔中央强度之比。（3）求$a=2b$时第一暗环的角半径。

\begin{center}\begin{minipage}{.98\linewidth}\centering
\begin{tikzpicture}[scale=.7,transform shape]
\draw[object] (0.00000,0.00000) circle(1.50000);\draw[object] (0.00000,0.00000) circle(0.75000);\draw[thin] (0.00000,0.00000)--(1.06050,1.06050);\node[above left] at (0.75000,0.75000) {$a$};\draw[thin] (0.00000,0.00000)--(0.53025,-0.53025);\node[below right] at (0.45000,-0.45000) {$b$};\begin{scope}[shift={(4.20000,0.00000)}]
\draw[axis] (-2.30000,0.00000)--(7.00000,0.00000);\draw[object] (0.00000,-1.20000)--(0.00000,-1.80000);\draw[object] (0.00000,1.20000)--(0.00000,1.80000);\draw[object] (2.80000,-2.40000)--(2.80000,2.40000);\draw[object] (2.67000,-2.20000)--(2.80000,-2.40000)--(2.93000,-2.20000);\draw[object] (2.67000,2.20000)--(2.80000,2.40000)--(2.93000,2.20000);\node[above] at (2.80000,2.50000) {$L$};\draw[object] (6.00000,-2.00000)--(6.00000,2.00000);\node[above] at (6.00000,2.00000) {焦平面};\fill (6.00000,0.00000) circle(.035);\node[below right] at (6.00000,0.00000) {$O$};\fill (6.00000,1.20000) circle(.035);\node[right] at (6.00000,1.20000) {$P$};\draw[very thick,color=muted] (0.00000,-0.60000)--(0.00000,0.60000);\draw[motion] (-2.00000,-1.00000)--(0.00000,-1.00000);
\draw[motion] (0.00000,-1.00000)--(2.80000,0.05000);
\draw[motion] (2.80000,0.05000)--(6.00000,1.20000);\draw[motion] (-2.00000,1.00000)--(0.00000,1.00000);
\draw[motion] (0.00000,1.00000)--(2.80000,2.05000);
\draw[motion] (2.80000,2.05000)--(6.00000,1.20000);\draw[<->,thin,color=muted] (2.80000,-2.30000)--(6.00000,-2.30000);\node[below] at (4.40000,-2.30000) {$f$};\draw[<->,thin,color=muted] (6.70000,0.00000)--(6.70000,1.20000);\node[right] at (6.70000,0.60000) {$r$};
\end{scope}
\end{tikzpicture}
\figcaption{图光3.10\quad 中央遮挡及同一外孔径的远场}
\end{minipage}\end{center}

\workarea

\clearpage

\pdfbookmark[2]{第 11 题 三种接收孔的孔径标度}{q11}

\section*{第 11 题\quad 三种接收孔的孔径标度}\label{q:11}

圆孔夫琅禾费实验中，衍射孔半径为$a$，入射强度、波长和透镜焦距固定。焦平面前放接收遮板，通光部分分别为：（1）中心很小的固定半径圆孔；（2）固定张角、半径足够大的扇形；（3）从中心沿半径伸出的固定窄宽长缝。证明接收功率在相应尺度条件下依次正比于$a^4,a^2,a^3$，并说明这些标度的适用范围。

\begin{center}\begin{minipage}{.98\linewidth}\centering
\begin{tikzpicture}[scale=.83,transform shape]
\begin{scope}[shift={(0.000000000,0.000000000)}]
\draw[object] (-1.400000000,-1.400000000)--(1.400000000,-1.400000000)--(1.400000000,1.400000000)--(-1.400000000,1.400000000)--cycle;
\draw[object] (0.000000000,0.000000000) circle(0.35);
\draw[guide] (0.000000000,0.000000000)--(0.350000000,0.000000000);
\draw[guide] (0.000000000,0.000000000)--(-0.550000000,-0.550000000);
\fill (0,0) circle(1pt);
\node[below left] at (-0.550000000,-0.550000000) {$O$};
\draw[<->,thin,color=muted] (0.000000000,-0.650000000)--(0.350000000,-0.650000000);
\node[below] at (0.175000000,-0.650000000) {$r_c$};
\node[] at (0.000000000,1.700000000) {中心小圆孔};
\end{scope}
\begin{scope}[shift={(4.200000000,0.000000000)}]
\draw[object] (-1.400000000,-1.400000000)--(1.400000000,-1.400000000)--(1.400000000,1.400000000)--(-1.400000000,1.400000000)--cycle;
\draw[object] (0,0)--(1.2,0) arc[start angle=0,end angle=55,radius=1.2]--cycle;
\node[] at (0.500000000,0.200000000) {$\varphi$};
\node[below left] at (0.000000000,0.000000000) {$O$};
\node[] at (0.000000000,1.700000000) {固定角扇形};
\end{scope}
\begin{scope}[shift={(8.400000000,0.000000000)}]
\draw[object] (-1.400000000,-1.400000000)--(1.400000000,-1.400000000)--(1.400000000,1.400000000)--(-1.400000000,1.400000000)--cycle;
\draw[object] (0.000000000,-0.120000000)--(1.200000000,-0.120000000)--(1.200000000,0.120000000)--(0.000000000,0.120000000)--cycle;
\node[below left] at (0.000000000,0.000000000) {$O$};
\draw[guide] (1.200000000,-0.120000000)--(1.650000000,-0.120000000);
\draw[guide] (1.200000000,0.120000000)--(1.650000000,0.120000000);
\draw[<->,thin,color=muted] (1.550000000,-0.120000000)--(1.550000000,0.120000000);
\node[right] at (1.720000000,0.000000000) {$w$};
\node[] at (0.000000000,1.700000000) {单侧径向窄条};
\end{scope}
\end{tikzpicture}
\figcaption{图光3.11\quad 三种固定尺寸的接收孔形状}
\end{minipage}\end{center}

\Needspace{6\baselineskip}

\workarea

\clearpage

\pdfbookmark[2]{第 12 题 有缺级要求的光栅设计}{q12}

\section*{第 12 题\quad 有缺级要求的光栅设计}\label{q:12}

设计正入射透射光栅，要求：（1）$\lambda=600\,\mathrm{nm}$的二级谱线衍射角不超过$30^\circ$，且在此前提下二级角色散尽量大；（2）第三级缺级；（3）二级光谱附近能分辨$\Delta\lambda=0.0200\,\mathrm{nm}$。确定光栅常数$d$、一种合适的缝宽$a$、所需最少总缝数$N$及有效宽度，并列出600纳米的可观察谱线级次。

\begin{center}\begin{minipage}{.98\linewidth}\centering
\begin{tikzpicture}[scale=.8,transform shape]
\draw[axis] (-2.50000,0.00000)--(6.80000,0.00000);\draw[object] (0.00000,-2.40000)--(0.00000,3.20000);\draw[object] (6.00000,-2.40000)--(6.00000,3.20000);\node[above] at (6.00000,3.20000) {焦平面};\draw[very thick,color=white] (0.00000,1.28333333)--(0.00000,1.51666667);\draw[thin] (-0.16000,1.28333333)--(0.16000,1.28333333);\draw[thin] (-0.16000,1.51666667)--(0.16000,1.51666667);\draw[motion] (-2.20000,1.40000)--(0.00000,1.40000);
\draw[motion] (0.00000,1.40000)--(3.00000,2.60000);
\draw[motion] (3.00000,2.60000)--(6.00000,1.20000);\node[above left] at (-0.12000,1.50000) {$1$};\draw[very thick,color=white] (0.00000,0.58333333)--(0.00000,0.81666667);\draw[thin] (-0.16000,0.58333333)--(0.16000,0.58333333);\draw[thin] (-0.16000,0.81666667)--(0.16000,0.81666667);\draw[motion] (-2.20000,0.70000)--(0.00000,0.70000);
\draw[motion] (0.00000,0.70000)--(3.00000,1.90000);
\draw[motion] (3.00000,1.90000)--(6.00000,1.20000);\node[above left] at (-0.12000,0.80000) {$2$};\draw[very thick,color=white] (0.00000,-0.11666667)--(0.00000,0.11666667);\draw[thin] (-0.16000,-0.11666667)--(0.16000,-0.11666667);\draw[thin] (-0.16000,0.11666667)--(0.16000,0.11666667);\draw[motion] (-2.20000,0.00000)--(0.00000,0.00000);
\draw[motion] (0.00000,0.00000)--(3.00000,1.20000);
\draw[motion] (3.00000,1.20000)--(6.00000,1.20000);\node[above left] at (-0.12000,0.10000) {$3$};\draw[very thick,color=white] (0.00000,-0.81666667)--(0.00000,-0.58333333);\draw[thin] (-0.16000,-0.81666667)--(0.16000,-0.81666667);\draw[thin] (-0.16000,-0.58333333)--(0.16000,-0.58333333);\draw[motion] (-2.20000,-0.70000)--(0.00000,-0.70000);
\draw[motion] (0.00000,-0.70000)--(3.00000,0.50000);
\draw[motion] (3.00000,0.50000)--(6.00000,1.20000);\node[above left] at (-0.12000,-0.60000) {$4$};\draw[very thick,color=white] (0.00000,-1.51666667)--(0.00000,-1.28333333);\draw[thin] (-0.16000,-1.51666667)--(0.16000,-1.51666667);\draw[thin] (-0.16000,-1.28333333)--(0.16000,-1.28333333);\draw[motion] (-2.20000,-1.40000)--(0.00000,-1.40000);
\draw[motion] (0.00000,-1.40000)--(3.00000,-0.20000);
\draw[motion] (3.00000,-0.20000)--(6.00000,1.20000);\node[above left] at (-0.12000,-1.30000) {$5$};\fill (6.00000,1.20000) circle(.035);\node[right] at (6.00000,1.20000) {$P$};\fill (6.00000,0.00000) circle(.035);\node[right] at (6.00000,0.00000) {$O$};\node[above] at (2.00000,0.35000) {$\theta$};\draw[object] (3.00000,-3.00000)--(3.00000,3.00000);\draw[object] (2.87000,-2.80000)--(3.00000,-3.00000)--(3.13000,-2.80000);\draw[object] (2.87000,2.80000)--(3.00000,3.00000)--(3.13000,2.80000);\node[above] at (3.00000,3.10000) {$L$};\node[above] at (1.20000,3.80000) {$d$：中心距，$a$：缝宽};
\end{tikzpicture}
\figcaption{图光3.12\quad 有限宽度的等间距透射光栅}
\end{minipage}\end{center}

\workarea

\clearpage

\pdfbookmark[2]{第 13 题 有限光栅的次极大位置与强度}{q13}

\section*{第 13 题\quad 有限光栅的次极大位置与强度}\label{q:13}

理想等间距$N$缝光栅的干涉因子为$F_N(\gamma)=[\sin(N\gamma)/\sin\gamma]^2$，$\gamma=\pi d\sin\theta/\lambda$。推导次极大的精确方程和近似位置，求最接近主极大的次峰与主峰强度之比，说明有限$N$和大$N$时的区别。先仅讨论干涉因子，单缝包络另行考虑。

\begin{center}\begin{minipage}{.98\linewidth}\centering
\begin{tikzpicture}[scale=.9,transform shape]
\draw[axis] (0.00000,0.00000)--(9.25000,0.00000);\draw[axis] (0.00000,0.00000)--(0.00000,2.65000);\node[right] at (9.25000,0.00000) {$\gamma/\pi$};\node[above right] at (0.00000,2.65000) {$F_5/25$};\draw[traj] plot coordinates {(0.00000,2.40000) (0.01800,2.39924) (0.03600,2.39697) (0.05400,2.39319) (0.07200,2.38790) (0.09000,2.38111) (0.10800,2.37283) (0.12600,2.36308) (0.14400,2.35187) (0.16200,2.33922) (0.18000,2.32513) (0.19800,2.30965) (0.21600,2.29278) (0.23400,2.27455) (0.25200,2.25499) (0.27000,2.23413) (0.28800,2.21200) (0.30600,2.18863) (0.32400,2.16405) (0.34200,2.13831) (0.36000,2.11143) (0.37800,2.08345) (0.39600,2.05443) (0.41400,2.02439) (0.43200,1.99338) (0.45000,1.96145) (0.46800,1.92863) (0.48600,1.89498) (0.50400,1.86055) (0.52200,1.82537) (0.54000,1.78951) (0.55800,1.75300) (0.57600,1.71591) (0.59400,1.67827) (0.61200,1.64015) (0.63000,1.60158) (0.64800,1.56263) (0.66600,1.52335) (0.68400,1.48378) (0.70200,1.44398) (0.72000,1.40400) (0.73800,1.36389) (0.75600,1.32371) (0.77400,1.28349) (0.79200,1.24330) (0.81000,1.20318) (0.82800,1.16318) (0.84600,1.12335) (0.86400,1.08373) (0.88200,1.04437) (0.90000,1.00533) (0.91800,0.96663) (0.93600,0.92832) (0.95400,0.89045) (0.97200,0.85306) (0.99000,0.81618) (1.00800,0.77985) (1.02600,0.74411) (1.04400,0.70899) (1.06200,0.67453) (1.08000,0.64076) (1.09800,0.60770) (1.11600,0.57540) (1.13400,0.54386) (1.15200,0.51312) (1.17000,0.48320) (1.18800,0.45412) (1.20600,0.42591) (1.22400,0.39857) (1.24200,0.37213) (1.26000,0.34659) (1.27800,0.32198) (1.29600,0.29829) (1.31400,0.27554) (1.33200,0.25374) (1.35000,0.23289) (1.36800,0.21299) (1.38600,0.19404) (1.40400,0.17605) (1.42200,0.15901) (1.44000,0.14291) (1.45800,0.12775) (1.47600,0.11353) (1.49400,0.10023) (1.51200,0.08784) (1.53000,0.07636) (1.54800,0.06576) (1.56600,0.05604) (1.58400,0.04717) (1.60200,0.03914) (1.62000,0.03193) (1.63800,0.02552) (1.65600,0.01989) (1.67400,0.01501) (1.69200,0.01087) (1.71000,0.00744) (1.72800,0.00469) (1.74600,0.00259) (1.76400,0.00113) (1.78200,0.00028) (1.80000,0.00000) (1.81800,0.00027) (1.83600,0.00106) (1.85400,0.00234) (1.87200,0.00408) (1.89000,0.00625) (1.90800,0.00883) (1.92600,0.01178) (1.94400,0.01507) (1.96200,0.01867) (1.98000,0.02256) (1.99800,0.02670) (2.01600,0.03107) (2.03400,0.03564) (2.05200,0.04037) (2.07000,0.04524) (2.08800,0.05023) (2.10600,0.05530) (2.12400,0.06043) (2.14200,0.06560) (2.16000,0.07078) (2.17800,0.07594) (2.19600,0.08107) (2.21400,0.08613) (2.23200,0.09112) (2.25000,0.09600) (2.26800,0.10076) (2.28600,0.10538) (2.30400,0.10985) (2.32200,0.11414) (2.34000,0.11824) (2.35800,0.12214) (2.37600,0.12582) (2.39400,0.12927) (2.41200,0.13249) (2.43000,0.13545) (2.44800,0.13816) (2.46600,0.14060) (2.48400,0.14276) (2.50200,0.14465) (2.52000,0.14626) (2.53800,0.14758) (2.55600,0.14862) (2.57400,0.14937) (2.59200,0.14983) (2.61000,0.15000) (2.62800,0.14989) (2.64600,0.14949) (2.66400,0.14883) (2.68200,0.14788) (2.70000,0.14667) (2.71800,0.14521) (2.73600,0.14349) (2.75400,0.14152) (2.77200,0.13933) (2.79000,0.13690) (2.80800,0.13427) (2.82600,0.13143) (2.84400,0.12839) (2.86200,0.12518) (2.88000,0.12180) (2.89800,0.11827) (2.91600,0.11460) (2.93400,0.11080) (2.95200,0.10688) (2.97000,0.10287) (2.98800,0.09877) (3.00600,0.09460) (3.02400,0.09038) (3.04200,0.08611) (3.06000,0.08182) (3.07800,0.07752) (3.09600,0.07322) (3.11400,0.06893) (3.13200,0.06467) (3.15000,0.06046) (3.16800,0.05631) (3.18600,0.05222) (3.20400,0.04822) (3.22200,0.04431) (3.24000,0.04051) (3.25800,0.03683) (3.27600,0.03328) (3.29400,0.02986) (3.31200,0.02660) (3.33000,0.02349) (3.34800,0.02055) (3.36600,0.01779) (3.38400,0.01520) (3.40200,0.01281) (3.42000,0.01060) (3.43800,0.00860) (3.45600,0.00680) (3.47400,0.00521) (3.49200,0.00383) (3.51000,0.00265) (3.52800,0.00170) (3.54600,0.00095) (3.56400,0.00042) (3.58200,0.00011) (3.60000,0.00000) (3.61800,0.00010) (3.63600,0.00042) (3.65400,0.00093) (3.67200,0.00164) (3.69000,0.00255) (3.70800,0.00364) (3.72600,0.00492) (3.74400,0.00637) (3.76200,0.00799) (3.78000,0.00977) (3.79800,0.01170) (3.81600,0.01378) (3.83400,0.01599) (3.85200,0.01833) (3.87000,0.02077) (3.88800,0.02333) (3.90600,0.02598) (3.92400,0.02871) (3.94200,0.03151) (3.96000,0.03437) (3.97800,0.03729) (3.99600,0.04024) (4.01400,0.04322) (4.03200,0.04621) (4.05000,0.04920) (4.06800,0.05219) (4.08600,0.05516) (4.10400,0.05810) (4.12200,0.06099) (4.14000,0.06384) (4.15800,0.06661) (4.17600,0.06932) (4.19400,0.07194) (4.21200,0.07447) (4.23000,0.07689) (4.24800,0.07921) (4.26600,0.08140) (4.28400,0.08346) (4.30200,0.08539) (4.32000,0.08718) (4.33800,0.08881) (4.35600,0.09029) (4.37400,0.09161) (4.39200,0.09276) (4.41000,0.09374) (4.42800,0.09455) (4.44600,0.09518) (4.46400,0.09564) (4.48200,0.09591) (4.50000,0.09600) (4.51800,0.09591) (4.53600,0.09564) (4.55400,0.09518) (4.57200,0.09455) (4.59000,0.09374) (4.60800,0.09276) (4.62600,0.09161) (4.64400,0.09029) (4.66200,0.08881) (4.68000,0.08718) (4.69800,0.08539) (4.71600,0.08346) (4.73400,0.08140) (4.75200,0.07921) (4.77000,0.07689) (4.78800,0.07447) (4.80600,0.07194) (4.82400,0.06932) (4.84200,0.06661) (4.86000,0.06384) (4.87800,0.06099) (4.89600,0.05810) (4.91400,0.05516) (4.93200,0.05219) (4.95000,0.04920) (4.96800,0.04621) (4.98600,0.04322) (5.00400,0.04024) (5.02200,0.03729) (5.04000,0.03437) (5.05800,0.03151) (5.07600,0.02871) (5.09400,0.02598) (5.11200,0.02333) (5.13000,0.02077) (5.14800,0.01833) (5.16600,0.01599) (5.18400,0.01378) (5.20200,0.01170) (5.22000,0.00977) (5.23800,0.00799) (5.25600,0.00637) (5.27400,0.00492) (5.29200,0.00364) (5.31000,0.00255) (5.32800,0.00164) (5.34600,0.00093) (5.36400,0.00042) (5.38200,0.00010) (5.40000,0.00000) (5.41800,0.00011) (5.43600,0.00042) (5.45400,0.00095) (5.47200,0.00170) (5.49000,0.00265) (5.50800,0.00383) (5.52600,0.00521) (5.54400,0.00680) (5.56200,0.00860) (5.58000,0.01060) (5.59800,0.01281) (5.61600,0.01520) (5.63400,0.01779) (5.65200,0.02055) (5.67000,0.02349) (5.68800,0.02660) (5.70600,0.02986) (5.72400,0.03328) (5.74200,0.03683) (5.76000,0.04051) (5.77800,0.04431) (5.79600,0.04822) (5.81400,0.05222) (5.83200,0.05631) (5.85000,0.06046) (5.86800,0.06467) (5.88600,0.06893) (5.90400,0.07322) (5.92200,0.07752) (5.94000,0.08182) (5.95800,0.08611) (5.97600,0.09038) (5.99400,0.09460) (6.01200,0.09877) (6.03000,0.10287) (6.04800,0.10688) (6.06600,0.11080) (6.08400,0.11460) (6.10200,0.11827) (6.12000,0.12180) (6.13800,0.12518) (6.15600,0.12839) (6.17400,0.13143) (6.19200,0.13427) (6.21000,0.13690) (6.22800,0.13933) (6.24600,0.14152) (6.26400,0.14349) (6.28200,0.14521) (6.30000,0.14667) (6.31800,0.14788) (6.33600,0.14883) (6.35400,0.14949) (6.37200,0.14989) (6.39000,0.15000) (6.40800,0.14983) (6.42600,0.14937) (6.44400,0.14862) (6.46200,0.14758) (6.48000,0.14626) (6.49800,0.14465) (6.51600,0.14276) (6.53400,0.14060) (6.55200,0.13816) (6.57000,0.13545) (6.58800,0.13249) (6.60600,0.12927) (6.62400,0.12582) (6.64200,0.12214) (6.66000,0.11824) (6.67800,0.11414) (6.69600,0.10985) (6.71400,0.10538) (6.73200,0.10076) (6.75000,0.09600) (6.76800,0.09112) (6.78600,0.08613) (6.80400,0.08107) (6.82200,0.07594) (6.84000,0.07078) (6.85800,0.06560) (6.87600,0.06043) (6.89400,0.05530) (6.91200,0.05023) (6.93000,0.04524) (6.94800,0.04037) (6.96600,0.03564) (6.98400,0.03107) (7.00200,0.02670) (7.02000,0.02256) (7.03800,0.01867) (7.05600,0.01507) (7.07400,0.01178) (7.09200,0.00883) (7.11000,0.00625) (7.12800,0.00408) (7.14600,0.00234) (7.16400,0.00106) (7.18200,0.00027) (7.20000,0.00000) (7.21800,0.00028) (7.23600,0.00113) (7.25400,0.00259) (7.27200,0.00469) (7.29000,0.00744) (7.30800,0.01087) (7.32600,0.01501) (7.34400,0.01989) (7.36200,0.02552) (7.38000,0.03193) (7.39800,0.03914) (7.41600,0.04717) (7.43400,0.05604) (7.45200,0.06576) (7.47000,0.07636) (7.48800,0.08784) (7.50600,0.10023) (7.52400,0.11353) (7.54200,0.12775) (7.56000,0.14291) (7.57800,0.15901) (7.59600,0.17605) (7.61400,0.19404) (7.63200,0.21299) (7.65000,0.23289) (7.66800,0.25374) (7.68600,0.27554) (7.70400,0.29829) (7.72200,0.32198) (7.74000,0.34659) (7.75800,0.37213) (7.77600,0.39857) (7.79400,0.42591) (7.81200,0.45412) (7.83000,0.48320) (7.84800,0.51312) (7.86600,0.54386) (7.88400,0.57540) (7.90200,0.60770) (7.92000,0.64076) (7.93800,0.67453) (7.95600,0.70899) (7.97400,0.74411) (7.99200,0.77985) (8.01000,0.81618) (8.02800,0.85306) (8.04600,0.89045) (8.06400,0.92832) (8.08200,0.96663) (8.10000,1.00533) (8.11800,1.04437) (8.13600,1.08373) (8.15400,1.12335) (8.17200,1.16318) (8.19000,1.20318) (8.20800,1.24330) (8.22600,1.28349) (8.24400,1.32371) (8.26200,1.36389) (8.28000,1.40400) (8.29800,1.44398) (8.31600,1.48378) (8.33400,1.52335) (8.35200,1.56263) (8.37000,1.60158) (8.38800,1.64015) (8.40600,1.67827) (8.42400,1.71591) (8.44200,1.75300) (8.46000,1.78951) (8.47800,1.82537) (8.49600,1.86055) (8.51400,1.89498) (8.53200,1.92863) (8.55000,1.96145) (8.56800,1.99338) (8.58600,2.02439) (8.60400,2.05443) (8.62200,2.08345) (8.64000,2.11143) (8.65800,2.13831) (8.67600,2.16405) (8.69400,2.18863) (8.71200,2.21200) (8.73000,2.23413) (8.74800,2.25499) (8.76600,2.27455) (8.78400,2.29278) (8.80200,2.30965) (8.82000,2.32513) (8.83800,2.33922) (8.85600,2.35187) (8.87400,2.36308) (8.89200,2.37283) (8.91000,2.38111) (8.92800,2.38790) (8.94600,2.39319) (8.96400,2.39697) (8.98200,2.39924) (9.00000,2.40000)};\draw[thin] (0.00000,-0.06000)--(0.00000,0.06000);\node[below] at (0.00000,-0.13000) {0};\draw[thin] (1.80000,-0.06000)--(1.80000,0.06000);\node[below] at (1.80000,-0.13000) {$1/5$};\draw[thin] (3.60000,-0.06000)--(3.60000,0.06000);\node[below] at (3.60000,-0.13000) {$2/5$};\draw[thin] (5.40000,-0.06000)--(5.40000,0.06000);\node[below] at (5.40000,-0.13000) {$3/5$};\draw[thin] (7.20000,-0.06000)--(7.20000,0.06000);\node[below] at (7.20000,-0.13000) {$4/5$};\draw[thin] (9.00000,-0.06000)--(9.00000,0.06000);\node[below] at (9.00000,-0.13000) {1};\node[above] at (4.50000,3.00000) {$N=5$：四个零点、三个次峰};
\end{tikzpicture}
\figcaption{图光3.13\quad 有限五缝干涉因子；尚未乘单缝包络}
\end{minipage}\end{center}

\workarea

\clearpage

\pdfbookmark[2]{第 14 题 显微镜的分辨距离与有效放大}{q14}

\section*{第 14 题\quad 显微镜的分辨距离与有效放大}\label{q:14}

空气中的显微镜物镜数值孔径$\mathrm{NA}=0.650$，照明波长$550\,\mathrm{nm}$。（1）按瑞利判据求最小分辨距离。（2）换用折射率$n=1.60$的油浸物镜，假定物方孔径半角的正弦仍为0.650，求分辨距离。（3）用空气物镜改以420纳米照明，求分辨距离。（4）要使人眼充分辨认第（2）问的细节，瞳孔直径2.50毫米、明视距离25.0厘米，按同样衍射判据估计所需最小视角放大率。

\begin{center}\begin{minipage}{.98\linewidth}\centering
\begin{tikzpicture}[scale=.9,transform shape]
\draw[axis] (-3.50000,0.00000)--(6.00000,0.00000);\fill (-3.00000,0.00000) circle(.035);\node[below] at (-3.00000,0.00000) {$Q$};\draw[object] (0.00000,-1.60000)--(0.00000,1.60000);\draw[object] (-0.13000,-1.40000)--(0.00000,-1.60000)--(0.13000,-1.40000);\draw[object] (-0.13000,1.40000)--(0.00000,1.60000)--(0.13000,1.40000);\node[above] at (0.00000,1.70000) {物镜};\fill (5.00000,0.00000) circle(.035);\node[below] at (5.00000,0.00000) {$Q'$};\draw[motion] (-3.00000,0.00000)--(0.00000,1.50000);
\draw[motion] (0.00000,1.50000)--(5.00000,0.00000);\draw[motion] (-3.00000,0.00000)--(0.00000,-1.50000);
\draw[motion] (0.00000,-1.50000)--(5.00000,0.00000);\node[above] at (-1.70000,0.05000) {$u$};\node[above] at (2.00000,0.40000) {$u'$ };\node[above] at (-2.00000,1.60000) {$n$};\node[above] at (3.00000,1.60000) {$n'$};
\end{tikzpicture}
\figcaption{图光3.14\quad 物方和像方孔径半角的定义}
\end{minipage}\end{center}

\workarea

\clearpage

\pdfbookmark[2]{第 15 题 交替间隔的成对光栅}{q15}

\section*{第 15 题\quad 交替间隔的成对光栅}\label{q:15}

光栅共有$2N$条缝宽均为$a$的狭缝，相邻缝间不透明部分的宽度依次为$a,3a,a,3a,\ldots$。波长$\lambda$的单色平行光正入射。求夫琅禾费强度分布，用两种分组方法推导，并讨论各因子决定的极大、极小及缺级。

\begin{center}\begin{minipage}{.98\linewidth}\centering
\begin{tikzpicture}[scale=.85,transform shape]
\draw[object] (-0.27500,-1.00000)--(0.27500,-1.00000)--(0.27500,1.00000)--(-0.27500,1.00000)--(-0.27500,-1.00000);\draw[object] (0.82500,-1.00000)--(1.37500,-1.00000)--(1.37500,1.00000)--(0.82500,1.00000)--(0.82500,-1.00000);\draw[object] (3.02500,-1.00000)--(3.57500,-1.00000)--(3.57500,1.00000)--(3.02500,1.00000)--(3.02500,-1.00000);\draw[object] (4.12500,-1.00000)--(4.67500,-1.00000)--(4.67500,1.00000)--(4.12500,1.00000)--(4.12500,-1.00000);\draw[object] (6.32500,-1.00000)--(6.87500,-1.00000)--(6.87500,1.00000)--(6.32500,1.00000)--(6.32500,-1.00000);\draw[object] (7.42500,-1.00000)--(7.97500,-1.00000)--(7.97500,1.00000)--(7.42500,1.00000)--(7.42500,-1.00000);\draw[<->,thin,color=muted] (0.00000,-1.50000)--(1.10000,-1.50000);\node[below] at (0.55000,-1.50000) {$2a$};\draw[<->,thin,color=muted] (1.10000,-1.50000)--(3.30000,-1.50000);\node[below] at (2.20000,-1.50000) {$4a$};\draw[<->,thin,color=muted] (0.00000,1.60000)--(3.30000,1.60000);\node[above] at (1.65000,1.60000) {$6a$};\node[above] at (4.80000,2.10000) {缝间不透明部分为$a,3a$交替};
\draw[<->,thin,color=accent] (-.275,0)--(.275,0);\node[above,font=\scriptsize] at (0,.05) {$a$};
\end{tikzpicture}
\figcaption{图光3.15\quad 六条缝示意交替间隙，实际总缝数为2N}
\end{minipage}\end{center}

\workarea

\clearpage

\pdfbookmark[2]{第 16 题 中部间距误差对光栅谱峰的影响}{q16}

\section*{第 16 题\quad 中部间距误差对光栅谱峰的影响}\label{q:16}

光栅原有$2N$条等宽缝、正常中心距$d$。中部第$N$条与第$N+1$条缝中心距变为$l$，两侧其余间距均为$d$。波长$\lambda$的单色平行光正入射。求强度分布，与完美光栅比较，分析主级强度、峰位、暗级和峰宽，说明是否可以一概认为分辨本领减半。

各缝宽度记为$a$，取$0<a\le d$且$l\ge a$，保证相邻透光缝不重叠。

\begin{center}\begin{minipage}{.98\linewidth}\centering
\begin{tikzpicture}[scale=.85,transform shape]
\draw[thick] (0.00000,-1.00000)--(0.00000,1.00000);\node[below] at (0.00000,-1.30000) {1};\draw[thick] (1.00000,-1.00000)--(1.00000,1.00000);\node[below] at (1.00000,-1.30000) {2};\draw[thick] (2.00000,-1.00000)--(2.00000,1.00000);\node[below] at (2.00000,-1.30000) {3};\draw[thick] (3.00000,-1.00000)--(3.00000,1.00000);\node[below] at (3.00000,-1.30000) {4};\draw[thick] (4.70000,-1.00000)--(4.70000,1.00000);\node[below] at (4.70000,-1.30000) {5};\draw[thick] (5.70000,-1.00000)--(5.70000,1.00000);\node[below] at (5.70000,-1.30000) {6};\draw[thick] (6.70000,-1.00000)--(6.70000,1.00000);\node[below] at (6.70000,-1.30000) {7};\draw[thick] (7.70000,-1.00000)--(7.70000,1.00000);\node[below] at (7.70000,-1.30000) {8};\draw[<->,thin,color=muted] (3.00000,1.50000)--(4.70000,1.50000);\node[above] at (3.85000,1.50000) {$l$};\draw[<->,thin,color=muted] (0.00000,1.50000)--(1.00000,1.50000);\node[above] at (0.50000,1.50000) {$d$};\draw[<->,thin,color=muted] (1.50000,-2.40000)--(6.20000,-2.40000);\node[below] at (3.85000,-2.40000) {$L=(N-1)d+l$};\node[above] at (4.00000,2.50000) {两组各N条，图取N等于4};
\end{tikzpicture}
\figcaption{图光3.16\quad 仅中部中心距改变的两组光栅}
\end{minipage}\end{center}

\workarea

\clearpage

\pdfbookmark[2]{第 17 题 两种波长的谱线重叠}{q17}

\section*{第 17 题\quad 两种波长的谱线重叠}\label{q:17}

一透射光栅每毫米有300条缝。正入射光含红光与一束短波光，在$24.46^\circ$方向看到二者谱线重叠。（1）按该重叠为红光二级与短波光三级，求两波长及其他非零重叠角。（2）求仅含红光的谱线方向。并说明如果只给“红光、紫光”两个颜色名称，级次判断存在哪些限制。

\begin{center}\begin{minipage}{.98\linewidth}\centering
\begin{tikzpicture}[scale=.82,transform shape]
\draw[axis] (-2.5,0)--(6.9,0) node[right] {法线};
\draw[object] (0,-2.2)--(0,3.3) node[above] {光栅};
\draw[object] (3,-2.2)--(3,3.3) node[above] {$L$};
\draw[object] (6,-2.2)--(6,3.3) node[above] {后焦平面};
\draw[motion] (-2.2,-1.400000000)--(0,-1.400000000);
\draw[motion] (0,-1.400000000)--(3,-0.035349802);
\draw[motion] (3,-0.035349802)--(6,1.364650198);
\draw[motion] (-2.2,-0.700000000)--(0,-0.700000000);
\draw[motion] (0,-0.700000000)--(3,0.664650198);
\draw[motion] (3,0.664650198)--(6,1.364650198);
\draw[motion] (-2.2,0.000000000)--(0,0.000000000);
\draw[motion] (0,0.000000000)--(3,1.364650198);
\draw[motion] (3,1.364650198)--(6,1.364650198);
\draw[motion] (-2.2,0.700000000)--(0,0.700000000);
\draw[motion] (0,0.700000000)--(3,2.064650198);
\draw[motion] (3,2.064650198)--(6,1.364650198);
\draw[motion] (-2.2,1.400000000)--(0,1.400000000);
\draw[motion] (0,1.400000000)--(3,2.764650198);
\draw[motion] (3,2.764650198)--(6,1.364650198);
\fill (6,1.364650198) circle(.04) node[right] {$P$};
\fill (6,0) circle(.035) node[below right] {$O$};
\draw[muted] (.9,0) arc[start angle=0,end angle=24.46,radius=.9];
\node at (1.65,.24) {$24.46^\circ$};
\node[above] at (-1.2,2.4) {正入射};
\node[below] at (2,-2.3) {两波长在同一方向的主极大};
\end{tikzpicture}
\figcaption{图光3.17\quad 同方向的平行衍射光经透镜会聚于$P$}
\end{minipage}\end{center}

\workarea

\clearpage

\pdfbookmark[2]{第 18 题 闪耀光栅的几何与角色散}{q18}

\section*{第 18 题\quad 闪耀光栅的几何与角色散}\label{q:18}

一闪耀反射光栅沿光栅平面的有效宽度为200毫米，每毫米500个刻槽，闪耀角$\theta_B=43^\circ26'$。（1）入射光垂直刻槽的工作槽面，求550纳米光在闪耀方向的级次及理想色分辨本领。（2）求该方向角色散。（3）若改为垂直光栅平均平面入射，求三级光谱中与槽面镜面反射方向重合的最强波长。忽略槽形及反射率随波长的细节。

\begin{center}\begin{minipage}{.98\linewidth}\centering
\begin{tikzpicture}[scale=.85,transform shape]
\draw[dashed,color=muted] (-3.00000,0.00000)--(3.00000,0.00000);\node[below] at (-2.80000,0.00000) {平均平面};\draw[very thick] (-1.08926,1.03127)--(1.08926,-1.03127);\draw[axis] (0.00000,0.00000)--(0.00000,3.00000);\node[above] at (0.00000,3.00000) {$\boldsymbol n_g$};\draw[axis] (0.00000,0.00000)--(2.06253,2.17852);\node[right] at (2.16566,2.28745) {$\boldsymbol n_f$};\draw[thin] (0,1) arc[start angle=90,end angle=46.566667,radius=1];\node[above] at (0.50000,1.30000) {$\theta_B$};\node[below] at (-0.20000,-1.40000) {工作槽面};\draw[motion] (1.84222,1.69856)--(0.12345,-0.11688);\draw[motion] (-0.12345,0.11688)--(1.59533,1.93231);\node[right] at (2.30000,1.50000) {入射、回返};
\end{tikzpicture}
\figcaption{图光3.18\quad 槽面法线与平均法线；光线在相邻槽面点上分开示意}
\end{minipage}\end{center}

\workarea

\clearpage

\pdfbookmark[2]{第 19 题 透射阶梯光栅的高阶光谱}{q19}

\section*{第 19 题\quad 透射阶梯光栅的高阶光谱}\label{q:19}

透射阶梯光栅由一系列厚度增量为$t$、折射率$n$的玻璃层构成，相邻台阶在横向伸出宽度$d$，共$N$阶。平行光正入射，用透镜观察夫琅禾费图样。衍射角很小。（1）推导强度分布和干涉主级方向。（2）取$n=1.50,t=0.500\,\mathrm{cm},\lambda=500\,\mathrm{nm},N=40$，求轴附近级次、角色散及色分辨本领。（3）说明其特点与近似条件。

\begin{center}\begin{minipage}{.98\linewidth}\centering
\begin{tikzpicture}[scale=.8,transform shape]
\draw[axis] (-2.20000,0.00000)--(6.20000,0.00000);\draw[object] (0.00000,-0.30000)--(0.65000,-0.30000)--(0.65000,0.30000)--(0.00000,0.30000)--(0.00000,-0.30000);\draw[motion] (-2.00000,0.00000)--(0.00000,0.00000);
\draw[motion] (0.00000,0.00000)--(0.65000,0.00000);
\draw[motion] (0.65000,0.00000)--(5.00000,1.08750);\draw[object] (0.00000,-1.00000)--(1.30000,-1.00000)--(1.30000,-0.40000)--(0.00000,-0.40000)--(0.00000,-1.00000);\draw[motion] (-2.00000,-0.70000)--(0.00000,-0.70000);
\draw[motion] (0.00000,-0.70000)--(1.30000,-0.70000);
\draw[motion] (1.30000,-0.70000)--(5.00000,0.22500);\draw[object] (0.00000,-1.70000)--(1.95000,-1.70000)--(1.95000,-1.10000)--(0.00000,-1.10000)--(0.00000,-1.70000);\draw[motion] (-2.00000,-1.40000)--(0.00000,-1.40000);
\draw[motion] (0.00000,-1.40000)--(1.95000,-1.40000);
\draw[motion] (1.95000,-1.40000)--(5.00000,-0.63750);\draw[object] (0.00000,-2.40000)--(2.60000,-2.40000)--(2.60000,-1.80000)--(0.00000,-1.80000)--(0.00000,-2.40000);\draw[motion] (-2.00000,-2.10000)--(0.00000,-2.10000);
\draw[motion] (0.00000,-2.10000)--(2.60000,-2.10000);
\draw[motion] (2.60000,-2.10000)--(5.00000,-1.50000);\node[above] at (1.20000,1.80000) {逐阶多出玻璃厚度 $t$};\node[above] at (4.60000,1.40000) {$\theta$};\draw[<->,thin,color=muted] (0.65000,-2.90000)--(1.30000,-2.90000);\node[below] at (0.97500,-2.90000) {$t$};\draw[<->,thin,color=muted] (3.00000,-0.70000)--(3.00000,0.00000);\node[right] at (3.00000,-0.35000) {$d$};
\end{tikzpicture}
\figcaption{图光3.19\quad 阶梯横向宽度与厚度增量；各阶只画工作带}
\end{minipage}\end{center}

\workarea

\clearpage

\pdfbookmark[2]{第 20 题 双平面波记录的全息光栅}{q20}

\section*{第 20 题\quad 双平面波记录的全息光栅}\label{q:20}

两束同频相干平行光，在感光片法线两侧、同一入射面内以角度$\theta_1,\theta_2$照明，振幅分别为$A,B$，波长为$\lambda$。假设显影后振幅透射率在工作范围内与曝光强度成线性关系。（1）求感光片上的振幅透射率。（2）取有效横向宽度$W=Nd$，以同波长平行光垂直再照明，求夫琅禾费强度分布，并说明可出现的衍射级与传播条件。

\begin{center}\begin{minipage}{.98\linewidth}\centering
\begin{tikzpicture}[scale=.85,transform shape]
\draw[axis] (-3.50000,0.00000)--(4.50000,0.00000);\draw[object] (0.00000,-2.00000)--(0.00000,2.00000);\node[above] at (0.00000,2.00000) {片面};\node[right] at (4.50000,0.00000) {$z$};\node[above] at (0.00000,2.50000) {$x$};\draw[motion] (-3.00000,-1.40000)--(0.00000,0.00000);\draw[motion] (-3.00000,1.10000)--(0.00000,0.00000);\draw[thin,color=muted] (-1,0) arc[start angle=180,end angle=205.0169,radius=1];\draw[thin,color=muted] (-1.3,0) arc[start angle=180,end angle=159.8637,radius=1.3];\node[below left] at (-1.05000,-0.12000) {$\theta_1$};\node[above left] at (-1.35000,0.12000) {$\theta_2$};\node[left] at (-3.20000,-1.40000) {$A$};\node[left] at (-3.20000,1.10000) {$B$};
\end{tikzpicture}
\figcaption{图光3.20\quad 两个片面投影相反的入射平面波}
\end{minipage}\end{center}

\workarea

\clearpage

\pdfbookmark[2]{第 21 题 全息光栅的反衬度与衍射效率}{q21}

\section*{第 21 题\quad 全息光栅的反衬度与衍射效率}\label{q:21}

双平面波记录全息光栅时两束强度比为$R=A^2/B^2$。显影后振幅透射率与曝光强度成线性关系$t=t_b+\eta I_{\rm exp}$。（1）求记录条纹的反衬度$\gamma$。（2）求每个一级的衍射效率，并说明在何种控制条件下它正比于$\gamma^2$。（3）结合感光材料的线性工作区，讨论束强比的选择，不把材料经验值当成普遍定律。

\begin{center}\begin{minipage}{.98\linewidth}\centering
\begin{tikzpicture}[scale=.9,transform shape]
\draw[axis] (0.00000,0.00000)--(9.25000,0.00000);\draw[axis] (0.00000,0.00000)--(0.00000,2.85000);\node[right] at (9.25000,0.00000) {$\ln R$};\node[above right] at (0.00000,2.85000) {$\gamma$};\draw[traj] plot coordinates {(0.00000,0.69109) (0.01800,0.69643) (0.03600,0.70182) (0.05400,0.70725) (0.07200,0.71271) (0.09000,0.71821) (0.10800,0.72376) (0.12600,0.72934) (0.14400,0.73496) (0.16200,0.74062) (0.18000,0.74632) (0.19800,0.75205) (0.21600,0.75783) (0.23400,0.76365) (0.25200,0.76951) (0.27000,0.77541) (0.28800,0.78136) (0.30600,0.78734) (0.32400,0.79336) (0.34200,0.79943) (0.36000,0.80553) (0.37800,0.81168) (0.39600,0.81787) (0.41400,0.82410) (0.43200,0.83038) (0.45000,0.83669) (0.46800,0.84305) (0.48600,0.84945) (0.50400,0.85590) (0.52200,0.86238) (0.54000,0.86891) (0.55800,0.87549) (0.57600,0.88210) (0.59400,0.88876) (0.61200,0.89547) (0.63000,0.90221) (0.64800,0.90901) (0.66600,0.91584) (0.68400,0.92272) (0.70200,0.92964) (0.72000,0.93661) (0.73800,0.94362) (0.75600,0.95068) (0.77400,0.95778) (0.79200,0.96493) (0.81000,0.97212) (0.82800,0.97935) (0.84600,0.98663) (0.86400,0.99396) (0.88200,1.00133) (0.90000,1.00874) (0.91800,1.01620) (0.93600,1.02371) (0.95400,1.03126) (0.97200,1.03886) (0.99000,1.04650) (1.00800,1.05418) (1.02600,1.06192) (1.04400,1.06969) (1.06200,1.07752) (1.08000,1.08538) (1.09800,1.09330) (1.11600,1.10125) (1.13400,1.10926) (1.15200,1.11731) (1.17000,1.12540) (1.18800,1.13354) (1.20600,1.14172) (1.22400,1.14995) (1.24200,1.15822) (1.26000,1.16654) (1.27800,1.17490) (1.29600,1.18331) (1.31400,1.19176) (1.33200,1.20026) (1.35000,1.20880) (1.36800,1.21738) (1.38600,1.22601) (1.40400,1.23468) (1.42200,1.24339) (1.44000,1.25215) (1.45800,1.26095) (1.47600,1.26979) (1.49400,1.27868) (1.51200,1.28761) (1.53000,1.29658) (1.54800,1.30559) (1.56600,1.31464) (1.58400,1.32374) (1.60200,1.33287) (1.62000,1.34205) (1.63800,1.35126) (1.65600,1.36052) (1.67400,1.36981) (1.69200,1.37915) (1.71000,1.38852) (1.72800,1.39793) (1.74600,1.40738) (1.76400,1.41686) (1.78200,1.42639) (1.80000,1.43594) (1.81800,1.44554) (1.83600,1.45517) (1.85400,1.46483) (1.87200,1.47453) (1.89000,1.48426) (1.90800,1.49403) (1.92600,1.50383) (1.94400,1.51366) (1.96200,1.52352) (1.98000,1.53341) (1.99800,1.54333) (2.01600,1.55328) (2.03400,1.56326) (2.05200,1.57327) (2.07000,1.58330) (2.08800,1.59336) (2.10600,1.60345) (2.12400,1.61356) (2.14200,1.62369) (2.16000,1.63385) (2.17800,1.64403) (2.19600,1.65423) (2.21400,1.66445) (2.23200,1.67468) (2.25000,1.68494) (2.26800,1.69522) (2.28600,1.70551) (2.30400,1.71581) (2.32200,1.72613) (2.34000,1.73647) (2.35800,1.74681) (2.37600,1.75717) (2.39400,1.76753) (2.41200,1.77791) (2.43000,1.78829) (2.44800,1.79868) (2.46600,1.80907) (2.48400,1.81946) (2.50200,1.82986) (2.52000,1.84026) (2.53800,1.85066) (2.55600,1.86106) (2.57400,1.87146) (2.59200,1.88185) (2.61000,1.89223) (2.62800,1.90261) (2.64600,1.91298) (2.66400,1.92334) (2.68200,1.93369) (2.70000,1.94402) (2.71800,1.95434) (2.73600,1.96464) (2.75400,1.97493) (2.77200,1.98519) (2.79000,1.99544) (2.80800,2.00566) (2.82600,2.01586) (2.84400,2.02603) (2.86200,2.03617) (2.88000,2.04629) (2.89800,2.05637) (2.91600,2.06642) (2.93400,2.07644) (2.95200,2.08642) (2.97000,2.09636) (2.98800,2.10626) (3.00600,2.11611) (3.02400,2.12593) (3.04200,2.13570) (3.06000,2.14542) (3.07800,2.15509) (3.09600,2.16470) (3.11400,2.17427) (3.13200,2.18378) (3.15000,2.19323) (3.16800,2.20262) (3.18600,2.21196) (3.20400,2.22122) (3.22200,2.23043) (3.24000,2.23956) (3.25800,2.24863) (3.27600,2.25762) (3.29400,2.26655) (3.31200,2.27539) (3.33000,2.28416) (3.34800,2.29285) (3.36600,2.30146) (3.38400,2.30998) (3.40200,2.31842) (3.42000,2.32677) (3.43800,2.33503) (3.45600,2.34320) (3.47400,2.35128) (3.49200,2.35926) (3.51000,2.36714) (3.52800,2.37492) (3.54600,2.38260) (3.56400,2.39018) (3.58200,2.39765) (3.60000,2.40502) (3.61800,2.41227) (3.63600,2.41942) (3.65400,2.42645) (3.67200,2.43336) (3.69000,2.44016) (3.70800,2.44684) (3.72600,2.45340) (3.74400,2.45984) (3.76200,2.46615) (3.78000,2.47233) (3.79800,2.47839) (3.81600,2.48432) (3.83400,2.49011) (3.85200,2.49578) (3.87000,2.50131) (3.88800,2.50670) (3.90600,2.51195) (3.92400,2.51707) (3.94200,2.52204) (3.96000,2.52688) (3.97800,2.53156) (3.99600,2.53611) (4.01400,2.54050) (4.03200,2.54475) (4.05000,2.54885) (4.06800,2.55280) (4.08600,2.55660) (4.10400,2.56024) (4.12200,2.56374) (4.14000,2.56707) (4.15800,2.57025) (4.17600,2.57327) (4.19400,2.57614) (4.21200,2.57885) (4.23000,2.58139) (4.24800,2.58378) (4.26600,2.58600) (4.28400,2.58807) (4.30200,2.58997) (4.32000,2.59170) (4.33800,2.59328) (4.35600,2.59468) (4.37400,2.59593) (4.39200,2.59701) (4.41000,2.59792) (4.42800,2.59867) (4.44600,2.59925) (4.46400,2.59967) (4.48200,2.59992) (4.50000,2.60000) (4.51800,2.59992) (4.53600,2.59967) (4.55400,2.59925) (4.57200,2.59867) (4.59000,2.59792) (4.60800,2.59701) (4.62600,2.59593) (4.64400,2.59468) (4.66200,2.59328) (4.68000,2.59170) (4.69800,2.58997) (4.71600,2.58807) (4.73400,2.58600) (4.75200,2.58378) (4.77000,2.58139) (4.78800,2.57885) (4.80600,2.57614) (4.82400,2.57327) (4.84200,2.57025) (4.86000,2.56707) (4.87800,2.56374) (4.89600,2.56024) (4.91400,2.55660) (4.93200,2.55280) (4.95000,2.54885) (4.96800,2.54475) (4.98600,2.54050) (5.00400,2.53611) (5.02200,2.53156) (5.04000,2.52688) (5.05800,2.52204) (5.07600,2.51707) (5.09400,2.51195) (5.11200,2.50670) (5.13000,2.50131) (5.14800,2.49578) (5.16600,2.49011) (5.18400,2.48432) (5.20200,2.47839) (5.22000,2.47233) (5.23800,2.46615) (5.25600,2.45984) (5.27400,2.45340) (5.29200,2.44684) (5.31000,2.44016) (5.32800,2.43336) (5.34600,2.42645) (5.36400,2.41942) (5.38200,2.41227) (5.40000,2.40502) (5.41800,2.39765) (5.43600,2.39018) (5.45400,2.38260) (5.47200,2.37492) (5.49000,2.36714) (5.50800,2.35926) (5.52600,2.35128) (5.54400,2.34320) (5.56200,2.33503) (5.58000,2.32677) (5.59800,2.31842) (5.61600,2.30998) (5.63400,2.30146) (5.65200,2.29285) (5.67000,2.28416) (5.68800,2.27539) (5.70600,2.26655) (5.72400,2.25762) (5.74200,2.24863) (5.76000,2.23956) (5.77800,2.23043) (5.79600,2.22122) (5.81400,2.21196) (5.83200,2.20262) (5.85000,2.19323) (5.86800,2.18378) (5.88600,2.17427) (5.90400,2.16470) (5.92200,2.15509) (5.94000,2.14542) (5.95800,2.13570) (5.97600,2.12593) (5.99400,2.11611) (6.01200,2.10626) (6.03000,2.09636) (6.04800,2.08642) (6.06600,2.07644) (6.08400,2.06642) (6.10200,2.05637) (6.12000,2.04629) (6.13800,2.03617) (6.15600,2.02603) (6.17400,2.01586) (6.19200,2.00566) (6.21000,1.99544) (6.22800,1.98519) (6.24600,1.97493) (6.26400,1.96464) (6.28200,1.95434) (6.30000,1.94402) (6.31800,1.93369) (6.33600,1.92334) (6.35400,1.91298) (6.37200,1.90261) (6.39000,1.89223) (6.40800,1.88185) (6.42600,1.87146) (6.44400,1.86106) (6.46200,1.85066) (6.48000,1.84026) (6.49800,1.82986) (6.51600,1.81946) (6.53400,1.80907) (6.55200,1.79868) (6.57000,1.78829) (6.58800,1.77791) (6.60600,1.76753) (6.62400,1.75717) (6.64200,1.74681) (6.66000,1.73647) (6.67800,1.72613) (6.69600,1.71581) (6.71400,1.70551) (6.73200,1.69522) (6.75000,1.68494) (6.76800,1.67468) (6.78600,1.66445) (6.80400,1.65423) (6.82200,1.64403) (6.84000,1.63385) (6.85800,1.62369) (6.87600,1.61356) (6.89400,1.60345) (6.91200,1.59336) (6.93000,1.58330) (6.94800,1.57327) (6.96600,1.56326) (6.98400,1.55328) (7.00200,1.54333) (7.02000,1.53341) (7.03800,1.52352) (7.05600,1.51366) (7.07400,1.50383) (7.09200,1.49403) (7.11000,1.48426) (7.12800,1.47453) (7.14600,1.46483) (7.16400,1.45517) (7.18200,1.44554) (7.20000,1.43594) (7.21800,1.42639) (7.23600,1.41686) (7.25400,1.40738) (7.27200,1.39793) (7.29000,1.38852) (7.30800,1.37915) (7.32600,1.36981) (7.34400,1.36052) (7.36200,1.35126) (7.38000,1.34205) (7.39800,1.33287) (7.41600,1.32374) (7.43400,1.31464) (7.45200,1.30559) (7.47000,1.29658) (7.48800,1.28761) (7.50600,1.27868) (7.52400,1.26979) (7.54200,1.26095) (7.56000,1.25215) (7.57800,1.24339) (7.59600,1.23468) (7.61400,1.22601) (7.63200,1.21738) (7.65000,1.20880) (7.66800,1.20026) (7.68600,1.19176) (7.70400,1.18331) (7.72200,1.17490) (7.74000,1.16654) (7.75800,1.15822) (7.77600,1.14995) (7.79400,1.14172) (7.81200,1.13354) (7.83000,1.12540) (7.84800,1.11731) (7.86600,1.10926) (7.88400,1.10125) (7.90200,1.09330) (7.92000,1.08538) (7.93800,1.07752) (7.95600,1.06969) (7.97400,1.06192) (7.99200,1.05418) (8.01000,1.04650) (8.02800,1.03886) (8.04600,1.03126) (8.06400,1.02371) (8.08200,1.01620) (8.10000,1.00874) (8.11800,1.00133) (8.13600,0.99396) (8.15400,0.98663) (8.17200,0.97935) (8.19000,0.97212) (8.20800,0.96493) (8.22600,0.95778) (8.24400,0.95068) (8.26200,0.94362) (8.28000,0.93661) (8.29800,0.92964) (8.31600,0.92272) (8.33400,0.91584) (8.35200,0.90901) (8.37000,0.90221) (8.38800,0.89547) (8.40600,0.88876) (8.42400,0.88210) (8.44200,0.87549) (8.46000,0.86891) (8.47800,0.86238) (8.49600,0.85590) (8.51400,0.84945) (8.53200,0.84305) (8.55000,0.83669) (8.56800,0.83038) (8.58600,0.82410) (8.60400,0.81787) (8.62200,0.81168) (8.64000,0.80553) (8.65800,0.79943) (8.67600,0.79336) (8.69400,0.78734) (8.71200,0.78136) (8.73000,0.77541) (8.74800,0.76951) (8.76600,0.76365) (8.78400,0.75783) (8.80200,0.75205) (8.82000,0.74632) (8.83800,0.74062) (8.85600,0.73496) (8.87400,0.72934) (8.89200,0.72376) (8.91000,0.71821) (8.92800,0.71271) (8.94600,0.70725) (8.96400,0.70182) (8.98200,0.69643) (9.00000,0.69109)};\draw[thin] (0.00000,-0.06000)--(0.00000,0.06000);\node[below] at (0.00000,-0.13000) {$-4$};\draw[thin] (4.50000,-0.06000)--(4.50000,0.06000);\node[below] at (4.50000,-0.13000) {0};\draw[thin] (9.00000,-0.06000)--(9.00000,0.06000);\node[below] at (9.00000,-0.13000) {4};
\end{tikzpicture}
\figcaption{图光3.21\quad 记录反衬度对束强比的对称依赖}
\end{minipage}\end{center}

\workarea

\clearpage

\pdfbookmark[2]{第 22 题 点物全息图的虚像与实像}{q22}

\section*{第 22 题\quad 点物全息图的虚像与实像}\label{q:22}

感光片放在$z=0$的$xy$平面，物点$Q=(x_0,y_0,-d_0)$，$d_0>0$。平行参考光在$xz$平面传播，其片面复振幅为$R(x)=A\exp(-ikx\sin\theta)$，$k=2\pi/\lambda$。物点发出的球面波与参考波共同曝光，显影后振幅透射率与曝光强度线性相关。以相同参考光再照明，在傍轴近似下证明一阶项形成虚像和实像，求位置并说明近似条件。

\begin{center}\begin{minipage}{.98\linewidth}\centering
\begin{tikzpicture}[scale=.8,transform shape]
\draw[axis] (-4.00000,0.00000)--(4.00000,0.00000);\draw[object] (0.00000,-2.00000)--(0.00000,2.40000);\node[right] at (4.00000,0.00000) {$z$};\node[above] at (0.00000,2.50000) {$x$};\fill (-3.00000,-0.80000) circle(.035);\node[below left] at (-3.00000,-0.80000) {$Q=(x_0,y_0,-d_0)$};\draw[motion] (-3.00000,-0.80000)--(0.00000,1.20000);\draw[motion] (-3.00000,-0.80000)--(0.00000,-1.20000);\draw[motion] (-3.00000,2.20000)--(0.00000,1.00000);\node[above] at (-2.00000,2.25000) {参考光 $R$};\draw[<->,thin,color=muted] (-3.00000,-2.20000)--(0.00000,-2.20000);\node[below] at (-1.50000,-2.20000) {$d_0$};\node[right] at (0.20000,2.00000) {全息片};
\end{tikzpicture}
\figcaption{图光3.22\quad 片面坐标、点物发散波及倾斜参考波}
\end{minipage}\end{center}

\workarea

\end{document}
